\documentclass[final,twocolumn,10pt]{elsarticle}

\usepackage[numbers]{natbib}
\usepackage[colorlinks=true, linkcolor=blue, citecolor=blue, urlcolor=blue]{hyperref}
\usepackage{rotating}
\usepackage{booktabs}
\usepackage{graphicx}
\usepackage{array}       % Cho định nghĩa cột với độ rộng cố định
\usepackage{makecell}    % Cho xuống dòng bằng \\ trong ô
\usepackage{adjustbox}   % Cho scale bảng
% \usepackage{subfigure}
\usepackage{subcaption}
\usepackage{algorithm}
\usepackage{algorithmic}
\usepackage{amsmath} 
\usepackage{cuted}
\usepackage{amsfonts}
\usepackage{cleveref}
\crefname{section}{Appendix}{Appendices}

% \usepackage[capitalize,noabbrev]{cleveref}
% \crefname{section}{Appendix}{Appendices}

\usepackage{amsmath, amssymb, amsthm}
\begin{document} % <-- Di chuyển dòng này lên trước
% \usepackage{amsthm} % hoặc \usepackage{ntheorem}
\newtheorem{definition}{Definition}


\begin{frontmatter}

\title{XXX}
\author{XXX}
\address{University of Information Technology, HCM City}

\begin{abstract}
% Anomaly detection in Internet of Things (IoT) networks is fundamentally constrained by the observability of anomalous behaviors in network traffic data. 
% In practice, IoT traffic exhibits heterogeneous structures, including global, cluster-level, and local deviations, while labels are often scarce or unreliable. 
% Under such conditions, anomaly detection claims are meaningful only if the assumed anomaly structures are reflected in observable signals.

% In this work, we propose the \textbf{Local One-Class Null Foley--Sammon Transformation (XXX)}, a one-class anomaly detection method that explicitly operates under the assumption that anomalies manifest as locally or globally separable structures in the feature space. 
% XXX operationalizes this assumption by constructing pseudo-classes via distance-based clustering, thereby exposing discriminative signals without relying on ground-truth anomaly labels. 
% A stable projection subspace is then obtained through Singular Value Decomposition (SVD), ensuring robustness in high-dimensional and potentially rank-deficient settings.

% By design, XXX does not aim to detect anomalies lacking measurable structural signals; rather, it provides a principled mechanism to exploit locally observable deviations when such signals are present. 
% The method employs efficient matrix operations with a computational complexity of $O(Nd^2)$, making it suitable for deployment in resource-constrained IoT environments.

% Empirical evaluations on four IoT intrusion detection datasets demonstrate that, under the stated assumptions, XXX achieves strong detection performance, with AUC-ROC and AUC-PR scores of up to $99.39\%$ and $99.52\%$, respectively. 
% These results indicate that when anomaly structures are observable and locally separable, the proposed approach yields robust detection performance even in the presence of annotation noise.
\end{abstract}


\begin{keyword}
Internet of Thing  \sep One-Class Novelty Detection \sep Null Foley-Sammon Transformation  \sep Network Intrusion Detection
\end{keyword}


\end{frontmatter}




\section{Introduction}

With more than 21 billion connected Internet of Things (IoT) devices projected worldwide by the end of 2025~\cite{iotanalytics2025devices}, large-scale IoT deployments continue to expand across domains such as healthcare, smart homes, and industrial automation. However, intrusion detection (IDS) in IoT networks remains difficult due to heterogeneous and resource-constrained devices. Machine learning (ML)-based IDS can model complex traffic patterns, but most approaches assume a closed-set setting and overfit known attack categories. 
 Consequently, detection
rates for unseen or zero-day attacks can drop by
up to 70\% in evaluation scenarios~\cite{Sarhan2023ZeroShot}. Meanwhile, recent generative technologies further accelerate the creation of exploit scripts and malware variants, increasing the frequency of previously unseen attacks~\cite{Ferrag2025GenAICybersec,Xu2025LLMAutonomousAttack}. This reveals a fundamental mismatch between closed-set IoT IDS assumptions and open-world realities, motivating a shift toward \emph{one-class novelty detection} (OCND) that models normal traffic and flags deviations as potential intrusions.

\subsection{One-Class Novelty Detection for IoT Intrusion Detection and Its Challenges}

Despite its conceptual suitability, reformulating IoT intrusion detection as OCND inevitably introduces a set of structural challenges in practical deployment:


\begin{itemize}
\item \textbf{Imperfect benign supervision.}  
In real-world deployments, benign training data are rarely perfectly clean due to labeling limitations and open-world uncertainty. Even slight contamination from asymmetric annotation errors can distort learned representations~\cite{Prabowo2026UncleanTrainingNIDS}.

\item \textbf{Structurally diverse anomaly manifestations.}  
Anomalies in IoT networks exhibit diverse structures due to heterogeneous, multi-source data, leading to significant variation in detector performance across anomaly types (Fig.~\ref{fig:ADBench_bar_chart})~\cite{han2022adbench,KANG2019772}.


\item \textbf{Constrained edge deployment regimes.}  
IoT IDS are often deployed at the network edge, where limited memory, computation, and latency budgets favor lightweight methods over computationally intensive deep architectures~\cite{Rahman2025IoTIDSSurvey,beggel2019robust,jodelka2021adaptive}.
\end{itemize}



Existing anomaly-detection methods operate under inherently conflicting design requirements in the OCND setting. Approaches that enforce compact representations  to improve sensitivity to subtle deviations are often highly vulnerable to slight contamination in the benign training set~\cite{perera2019ocgan,gardner2006one,han2022adbench}. In contrast, methods that explicitly enhance robustness to noisy supervision typically rely on more complex optimization procedures or relaxed representations, which increase computational overhead and limit their suitability for resource-constrained edge deployment~\cite{beggel2019robust,jodelka2021adaptive,Rahman2025IoTIDSSurvey}. Achieving robustness to benign-set contamination, adaptability to heterogeneous anomaly structures, and computational efficiency simultaneously remains fundamentally challenging, making practical OCND frameworks scarce in IoT IDS~\cite{perera2019ocgan,gardner2006one,han2022adbench,beggel2019robust,jodelka2021adaptive,Rahman2025IoTIDSSurvey} (see Section~\ref{sec:relatework}).


\begin{figure}[h]
    \centering
    \includegraphics[width=1\linewidth]{pictures/anomaly_model_comparison (1).png}
    \caption{Illustrative comparison of representative anomaly detectors under local, global, and clustered anomaly structures, adapted from the ADBench evaluation using average AUROC (\%). The figure is used to motivate that detector behavior can vary substantially across anomaly types. The best and second-best detector in each setting are highlighted in black and gray, respectively.}
    \label{fig:ADBench_bar_chart}
\end{figure}





\subsection{Proposed Approach}
We argue that practical IoT anomaly detection requires balancing two key trade-offs:
(i) maintaining lightweight training and inference under resource constraints, and
(ii) preserving local structural characteristics while remaining robust to noise and slight contamination in benign data. To address the efficiency requirement, we adopt a lightweight shadow learning paradigm that avoids costly model updates and enables fast deployment. Within this framework, prototype-based representations support efficient inference through simple distance computations to a small set of representatives. However, operating in the original feature space makes such representations highly sensitive to noise and limits their ability to capture intrinsic data structure.

To overcome this limitation, we introduce a Fisher-inspired null-space projection that suppresses intra-class variations while preserving underlying structural relationships. Specifically, we propose a \emph{linear localized discriminative modeling in the null space}, following the principle of the Null Foley--Sammon Transform (NFST), which maps diverse normal samples into a compact subspace with reduced irrelevant variations. Building on this foundation, \textbf{XXX} introduces two key adaptations:
\textbf{(1)} pseudo-class partitioning of benign data to recover latent discriminative structure without requiring labeled anomalies, and
\textbf{(2)} localized, representative-based evaluation in the projected space, enabling adaptive detection across diverse anomaly types.



However, directly extending null-space methods such as NFST to the OCND setting is non-trivial. A naive one-class formulation leads to \emph{null-space collapse}, where the lack of class contrast removes the discriminative structure required for meaningful projections, resulting in degenerate solutions. Moreover, IoT datasets often lie in the regime \(n \gg d\), where classical null-space construction becomes numerically unstable~\cite{ali2022theoretical,nguyen2025nnfst}. To address these issues, We locally partition benign traffic into pseudo-classes to recover the structural variability necessary for effective discriminative null-space learning, and adopt an SVD-based construction to ensure numerical stability in IoT environments. Together, this enables practical and robust null-space modeling for benign-only IoT intrusion detection.


\subsection{Main Contributions}

In summary, the main contributions of this work are as follows:


\begin{itemize}
    \item We reformulate open-world IoT intrusion detection from a closed-set classification paradigm to a benign-only one-class novelty detection (OCND) framework. This reformulation introduces three fundamental structural challenges which collectively define the practical design space for IoT IDS.

    \item We propose \textbf{XXX}, a lightweight OCND framework based on null-space projections. It combines pseudo-class construction, localized representative-based scoring, and SVD-based null-space computation. The resulting design is robust to slight benign-set contamination and heterogeneous anomaly behaviors, while maintaining low inference complexity for resource-constrained IoT environments.

    \item We evaluate XXX extensively on four standard IoT benchmarks (CICIoT2023, ToN-IoT, NBaIoT, and BoT-IoT), including controlled experiments with synthetic local, cluster, and global anomaly types and benign-set contamination rates of 1\%, 3\%, and 5\%. The results demonstrate that XXX achieves strong and consistent performance across all scenarios.
\end{itemize}

\section{Related Work}
\label{sec:relatework}

\begin{figure}[h]
    \centering
    \includegraphics[width=1\linewidth]{pictures/relatedWorks.pdf}
    \caption{Conceptual positioning of prior work related to anomaly detection for IoT intrusion detection. The literature is organized into general anomaly-detection families, benign-only novelty-detection methods, and robustness-oriented directions under non-ideal assumptions. The highlighted block indicates the position of our method.}
    \label{fig:relatedWorks}
\end{figure}

Intrusion Detection Systems (IDSs) are essential for securing IoT networks against threats such as distributed denial-of-service attacks, botnet activity, and probing. Due to the heterogeneous and partially labeled nature of IoT traffic, anomaly detection has emerged as a practical paradigm for identifying deviations from normal behavior~\cite{samara2022survey}. Compared with signature-based approaches, it is inherently better suited to open-world settings where attack patterns are incomplete or unseen.

\subsection{Anomaly Detection for IoT Intrusion Detection}

Existing approaches can be broadly divided into shallow and deep methods.  
Shallow models, including density-based (e.g., LOF), statistical (e.g., ECOD), and ensemble-based methods (e.g., Isolation Forest), are computationally efficient and easy to deploy~\cite{KNN,9338429,IForest}. However, their limited representational capacity often fails to capture the heterogeneous and multimodal structure of benign IoT traffic, leading to unstable performance.

Deep methods improve expressiveness through reconstruction-based models (e.g., Autoencoders), compactness-driven objectives (e.g., DeepSVDD), and adversarial frameworks (e.g., ALAD)~\cite{Autoencoder,ALAD}. While effective in modeling complex distributions, they introduce significant computational overhead and require careful tuning, limiting their practicality in resource-constrained environments. This reflects a fundamental trade-off between representation power and deployment efficiency.

\subsection{One-Class Novelty Detection and Null-Space Methods}

In open-world IoT security, intrusion detection is often formulated as a one-class novelty detection (OCND) problem, where models are trained on benign data only. Classical methods such as One-Class SVM estimate a decision boundary, while DeepSVDD learns a compact latent representation of normality~\cite{shin2005one,DeepSVDD}.

However, these approaches rely on a global notion of normality, implicitly assuming that benign data forms a single coherent structure. This assumption is often violated in IoT environments, where normal traffic spans multiple operating regimes. As a result, such models may either overfit to dominant patterns or miss locally structured anomalies.

Null-space-based methods provide an alternative perspective. Kernel NFST projects data into a discriminative null space by collapsing intra-class variance and measuring deviations in the projected space~\cite{bodesheim2013kernel}. Extensions such as nNFST further improve structure preservation for intrusion detection~\cite{nguyen2025nnfst}. Nevertheless, these methods are primarily designed for supervised or multiclass settings and do not naturally extend to benign-only learning, nor do they explicitly address contamination or scalability.

\subsection{Robust Detection under Non-Ideal Assumptions}

Recent studies have explored anomaly detection under more realistic conditions, including heterogeneous anomaly structures and imperfect supervision. Empirical evidence shows that anomalies may appear as global shifts, clustered deviations, or subtle local irregularities, and no single global scoring mechanism performs consistently across all cases~\cite{han2022adbench,KANG2019772}.

Another line of work focuses on robustness to contaminated nominal data. Methods such as DIF and DASVDD aim to mitigate sensitivity to noisy supervision~\cite{DIF,DASVDD}. However, these approaches typically introduce additional complexity or optimization overhead, limiting scalability and deployment in IoT scenarios~\cite{Rahman2025IoTIDSSurvey,beggel2019robust,jodelka2021adaptive}.

\textbf{Summary.} Existing methods address representation, robustness, and efficiency largely in isolation. There remains a lack of a unified framework that simultaneously captures heterogeneous benign structure, maintains robustness under slight contamination, and operates efficiently at scale. This gap motivates the proposed approach.

\section{Problem Formulation}

\subsection{Problem Definition: Benign-Only Novelty Detection under Practical Conditions}
\label{sec:problem_def}

Consider an IoT network monitoring environment with feature space $\mathcal{X}\subseteq\mathbb{R}^{d}$. Let $\mathbb{P}_{\text{norm}}$ denote the distribution of benign network traffic over $\mathcal{X}$. In the benign-only one-class novelty detection setting, we observe a training set
\[
\mathcal{D}_{\text{train}}=\{x_1,\ldots,x_N\}, \qquad x_i\in\mathbb{R}^{d},
\]
and seek to learn a scoring function
\[
S:\mathbb{R}^{d}\rightarrow\mathbb{R}
\]
such that benign samples receive low scores, while anomalous samples receive high scores.

Given a decision threshold $\tau$, a test sample $x_{\text{test}}$ is classified as
\begin{equation}
f(x_{\text{test}})=
\begin{cases}
0, & \text{if } S(x_{\text{test}})\le \tau,\\
1, & \text{if } S(x_{\text{test}})> \tau,
\end{cases}
\label{eq:ocnd_decision}
\end{equation}
where label $0$ denotes benign traffic and label $1$ denotes anomalous traffic. In practice, $\tau$ is typically chosen from the training-score distribution according to a target false-positive level or a fixed contamination assumption.

While this formulation is standard, practical IoT intrusion detection departs from the ideal OCND setting in three important ways.

\paragraph{1. Slight contamination in the benign training set.}
In real deployments, the benign training set may not be perfectly clean. Undetected attacks, mislabeled flows, or annotation bias can introduce a small number of anomalous samples into the nominal training data. A convenient way to express this is to model the training samples as being drawn from a contaminated distribution
\begin{equation}
\mathbb{P}_{\text{train}}=(1-\varepsilon)\mathbb{P}_{\text{norm}}+\varepsilon \mathbb{P}_{\text{cont}},
\qquad 0\le \varepsilon \ll 1,
\label{eq:contaminated_train}
\end{equation}
where $\mathbb{P}_{\text{cont}}$ denotes the contamination distribution. The resulting challenge is to estimate a useful benign support while reducing sensitivity to this slight contamination.

\paragraph{2. Multimodal benign structure and heterogeneous anomaly manifestations.}
Benign IoT traffic rarely follows a single homogeneous pattern because devices, protocols, and usage contexts can induce multiple operating regimes. A convenient abstraction is to model benign traffic as a mixture distribution
\begin{align}
\mathbb{P}_{\text{norm}}(x)
&= \sum_{k=1}^{K} \pi_k\, \mathbb{P}_k(x \mid \mu_k, \Sigma_k), \\
\pi_k &\ge 0,\quad \sum_{k=1}^{K} \pi_k = 1.
\label{eq:mixture_norm}
\end{align}
where each component $\mathbb{P}_k$ represents a latent benign traffic pattern with mean $\mu_k$ and covariance $\Sigma_k$.

Under such multimodal benign structure, anomalous behavior may appear in different forms, including local anomalies that deviate from the neighborhood structure of a benign mode, clustered anomalies that form coherent abnormal groups, and global anomalies that lie far from the overall benign support. Effective IoT novelty detection therefore requires more than a single global notion of deviation.

\paragraph{3. Deployment-oriented efficiency constraints.}
IoT intrusion detectors are often expected to operate under limited memory and computation budgets, especially in edge-oriented or large-scale monitoring environments. These constraints favor methods with efficient inference, compact model state, and practical memory usage on tabular traffic data. Thus, beyond detection performance alone, a practical OCND method should also remain lightweight enough for realistic deployment.

\paragraph{Problem statement.}
Given a training set drawn primarily from a multimodal benign distribution but subject to slight contamination, our goal is to learn a lightweight anomaly scoring function and decision rule that:
(i) preserves benign structure sufficiently well to support benign-only novelty detection,
(ii) remains empirically robust under slight training contamination, and
(iii) can respond to heterogeneous anomaly manifestations while maintaining practical inference efficiency on large tabular IoT traffic.

\subsection{Why Classical NFST is Nontrivial in Benign-Only Novelty Detection}
\label{sec:nfst_challenges}

The Null Foley--Sammon Transform (NFST) is a discriminative linear projection method derived from the Fisher criterion. Given a projection direction $w$, the Fisher objective is
\begin{equation}
F(w)=\frac{w^{\top}S_b w}{w^{\top}S_w w},
\label{eq:fisher}
\end{equation}
where $S_b$ and $S_w$ denote the between-class and within-class scatter matrices, respectively. For a multiclass training set with class means $\mu_j$, global mean $\mu$, and class size $n_j$, these matrices are
\begin{align}
S_b &= \sum_{j=1}^{c} n_j(\mu_j - \mu)(\mu_j - \mu)^{\top}, \\
S_w &= \sum_{j=1}^{c} \sum_{i=1}^{n_j} (x_{ij} - \mu_j)(x_{ij} - \mu_j)^{\top}.
\label{eq:scatter}
\end{align}

NFST seeks projection directions that annihilate the within-class scatter while preserving between-class separation:
\begin{equation}
w^{\top}S_w w = 0,
\qquad
w^{\top}S_b w > 0.
\label{eq:nfst_constraint}
\end{equation}
Directions satisfying Eq.~(\ref{eq:nfst_constraint}) are called \emph{null projecting directions} (NPDs).

Let $X_t=[x_1-\mu,\ldots,x_N-\mu]\in\mathbb{R}^{d\times N}$ be the centered data matrix, and let $Q\in\mathbb{R}^{d\times r}$ be an orthonormal basis for $\mathrm{span}(X_t)$. Any projection vector can be written as $w=Q\beta$, which leads to the reduced null-space condition
\begin{equation}
Q^{\top}S_w Q\,\beta = 0.
\label{eq:nullspace}
\end{equation}
Therefore, valid NPDs exist only if the null space $\mathcal{N}(Q^{\top}S_w Q)$ is nontrivial and contains directions that still preserve positive between-class scatter.

Although NFST is effective in multiclass discrimination, two difficulties arise when it is transferred to benign-only novelty detection.

\paragraph{1. Degeneracy under one-class training.}
Classical NFST fundamentally relies on between-class scatter. In benign-only training, however, there is only one observed class ($c=1$). In that case, the global mean coincides with the class mean, so $S_b=0$. Consequently, the discriminative condition $w^{\top}S_b w>0$ in Eq.~(\ref{eq:nfst_constraint}) cannot be satisfied, and the original NFST objective degenerates. This means that classical NFST cannot be applied directly without introducing additional structure.

\paragraph{2. Exact null directions may be unavailable or numerically fragile in large tabular settings.}
A second difficulty arises in large tabular regimes, especially when $N\gg d$. In such settings, the matrix $Q^{\top}S_w Q$ need not have a nontrivial exact null space; in particular, it may become full rank within the total-scatter subspace, leaving no valid exact NPDs. Even when near-null directions exist, their numerical recovery can be sensitive to basis construction and floating-point error. Therefore, the challenge is not only numerical instability in classical orthogonalization, but also the possibility that exact null-space directions may be absent or difficult to recover robustly in practice.

Taken together, these observations imply that adapting NFST to benign-only IoT novelty detection requires more than a direct reuse of the classical formulation. A practical reformulation must introduce surrogate discriminative structure for benign-only data and compute the projection in a numerically stable manner for large tabular datasets.




% ========================

\section{Proposed Methodology}
\label{sec:method}

This section presents \textbf{XXX}, a benign-only novelty-detection framework that reformulates the Null Foley--Sammon Transform (NFST) for tabular IoT intrusion detection. The method targets three requirements introduced in Section~\ref{sec:problem_def}: (i) applicability under benign-only training, (ii) accommodation of multimodal benign traffic under a lightweight scoring rule, and (iii) numerically stable projection learning in large tabular settings. 

For clarity, we present the method from two complementary perspectives. Figure~\ref{fig:workflow_system} shows a deployment-oriented view of an XXX-based IoT intrusion detection system, while Figure~\ref{fig:workflow_algorithm} summarizes the algorithmic workflow of XXX itself. As in the benchmark datasets used in this work, the methodological contribution of this paper begins from pre-extracted tabular traffic features rather than raw packets.

\begin{figure}[!h]
\centering
\includegraphics[width=\columnwidth]{pictures/Hihi.png}
\caption{Deployment-oriented workflow of the proposed XXX-based IoT intrusion detection system. Raw network traffic is first converted into flow- or session-level tabular features, after which the benign-only detector operates in the learned feature space.}
\label{fig:workflow_system}
\end{figure}

\subsection{Deployment-Oriented Workflow}
\label{subsec:deployment_workflow}

In practical deployment, XXX operates on flow-level or session-level IoT traffic features rather than raw packets. As illustrated in Figure~\ref{fig:workflow_system}, packet streams are captured at monitoring nodes, aggregated into flows or sessions, and transformed into fixed-dimensional feature vectors before being passed to the anomaly detector. While this reflects a complete deployment pipeline, the core methodological contribution of this work begins from the tabular feature representation, consistent with the benchmark datasets used in our experiments. 

From an operational viewpoint, the full system contains three stages:

\paragraph{1. Data preparation.}
IoT traffic is first represented as feature vectors containing statistical and protocol-level information, such as packet counts, byte counts, duration, or other flow descriptors. These features are then cleaned and normalized so that all samples lie in a common \(d\)-dimensional feature space.

\paragraph{2. Model training.}
Only benign training traffic is used to learn the detector. Since classical NFST requires multiple classes, XXX first partitions benign data into pseudo-classes and then learns a stable discriminative projection using an SVD-based NFST reformulation.

\paragraph{3. Online detection.}
For a new traffic sample, XXX computes a novelty score in the learned subspace by comparing the sample to projected representatives of benign pseudo-classes. The score is then compared with a decision threshold to determine whether the sample is normal or anomalous.

\subsection{Algorithmic Overview of XXX}
\label{subsec:overview_xxx}

Let the benign training set be
\[
\mathcal{D}_{\text{train}}=\{x_1,\ldots,x_n\}, \qquad x_i\in\mathbb{R}^{d},
\]
and let
\[
X=[x_1,\ldots,x_n]\in\mathbb{R}^{d\times n}
\]
denote the corresponding data matrix. XXX learns a projection matrix \(W\) and a novelty score \(s(\cdot)\) using only benign training data. The method consists of three stages:

\begin{enumerate}
    \item \textbf{Pseudo-class construction.} The benign training set is partitioned into \(c\ge 2\) pseudo-classes to restore a surrogate between-class structure for NFST.
    \item \textbf{SVD-based NFST reformulation.} A stable projection is computed inside the total-scatter subspace using exact or near-null directions of the reduced within-class scatter matrix.
    \item \textbf{Representative-based novelty scoring.} Test samples are scored by their distances to projected pseudo-class representatives, followed by a single global threshold estimated from benign training scores.
\end{enumerate}

Figure~\ref{fig:workflow_algorithm} summarizes this pipeline.

\begin{figure}[!ht]
    \centering
    \includegraphics[width=\columnwidth]{pictures/Workflowr.png}
    \caption{Algorithmic workflow of XXX. Benign training traffic is partitioned into pseudo-classes, scatter matrices are constructed from the pseudo-class structure, a stable SVD-based NFST projection is learned, and test samples are scored by their distances to projected pseudo-class representatives.}
    \label{fig:workflow_algorithm}
\end{figure}

Throughout the paper, \(c\) denotes the number of pseudo-classes, \(\epsilon\ge 0\) denotes the near-null tolerance used in the projection step, and \(\rho\in[0,1)\) denotes the upper-tail trimming parameter used for threshold estimation. Their selection protocol is specified in Section~\ref{sec:exp_su}.

\subsection{Pseudo-Class Construction}
\label{subsec:pseudoclass_xxx}

Classical NFST requires more than one class. To introduce surrogate class structure in the benign-only setting, XXX partitions the training set into \(c\ge 2\) clusters using K-means:
\[
\mathcal{D}_{\text{train}}=\bigcup_{j=1}^{c} C_j,
\qquad
C_j\cap C_\ell = \varnothing \text{ for } j\neq \ell,
\]
where \(C_j\) is the \(j\)-th pseudo-class and \(n_j=|C_j|\).

Let the pseudo-class centroid and global mean be
\[
\mu_j=\frac{1}{n_j}\sum_{x_i\in C_j}x_i,
\qquad
\mu=\frac{1}{n}\sum_{i=1}^{n}x_i.
\]

Figure~\ref{fig:pseudoclass_effect} illustrates the role of pseudo-class construction. Without clustering, benign data remain a single observed class and classical NFST degenerates. With pseudo-class partitioning, benign traffic is represented through multiple local regimes, which restores a surrogate between-class structure for projection learning and provides multiple benign representatives for later scoring.

\begin{figure}[!ht]
    \centering
    \begin{subfigure}{\columnwidth}
        \centering
        \includegraphics[width=\columnwidth]{pictures/Different1.png}
        \caption{Without pseudo-class construction, benign data remain a single observed class and classical NFST degenerates.}
        \label{fig:no_cluster_rewrite}
    \end{subfigure}

    \vspace{0.5em}

    \begin{subfigure}{\columnwidth}
        \centering
        \includegraphics[width=\columnwidth]{pictures/Different2.png}
        \caption{Pseudo-class partitioning introduces surrogate class structure and multiple benign modes for projection and scoring.}
        \label{fig:with_cluster_rewrite}
    \end{subfigure}

    \caption{Motivation for pseudo-class construction in XXX. Clustering benign traffic into \(c\) pseudo-classes avoids one-class degeneracy and yields multiple benign representatives in the projected space.}
    \label{fig:pseudoclass_effect}
\end{figure}

Using these pseudo-classes, XXX defines the within-class and between-class scatter matrices as
\begin{equation}
S_w=
\sum_{j=1}^{c}\sum_{x_i\in C_j}(x_i-\mu_j)(x_i-\mu_j)^\top,
\label{eq:sw_xxx_revised}
\end{equation}
and
\begin{equation}
S_b=
\sum_{j=1}^{c} n_j(\mu_j-\mu)(\mu_j-\mu)^\top.
\label{eq:sb_xxx_revised}
\end{equation}

Pseudo-class construction serves two purposes. First, it avoids the one-class degeneracy of classical NFST by restoring a nontrivial between-class scatter matrix. Second, it provides multiple benign representatives for subsequent scoring instead of relying on a single global benign prototype. We emphasize that K-means is used here as an efficient partitioning heuristic; it is not assumed to recover true latent classes or to guarantee robustness to arbitrary contamination patterns. Its effect is therefore assessed empirically in Section~\ref{sec:result_rq2}.

\subsection{Stable NFST Reformulation via Exact or Near-Null Directions}
\label{subsec:projection_xxx}

Let the centered data matrix be
\begin{equation}
X_c=[x_1-\mu,\ldots,x_n-\mu]\in\mathbb{R}^{d\times n}.
\label{eq:centered_matrix_xxx_revised}
\end{equation}
We compute its compact singular value decomposition
\begin{equation}
X_c=U_r\Sigma_r V_r^\top,
\label{eq:svd_xxx_revised}
\end{equation}
where \(U_r\in\mathbb{R}^{d\times r}\) is an orthonormal basis of the total-scatter subspace and \(r=\operatorname{rank}(X_c)\).

Any projection direction inside this subspace can be written as
\[
w=U_r\beta.
\]
Define the reduced within-class and between-class scatter matrices by
\begin{equation}
M=U_r^\top S_w U_r,
\qquad
A=U_r^\top S_b U_r.
\label{eq:reduced_scatter_xxx}
\end{equation}

In the classical NFST setting, valid directions are sought in the null space of \(M\). However, in large tabular regimes exact null directions may be unavailable or numerically fragile. To make the construction operational in both cases, XXX uses an exact-or-near-null subspace defined by the eigendecomposition
\begin{equation}
M=V\Lambda V^\top,
\qquad
\Lambda=\operatorname{diag}(\lambda_1,\ldots,\lambda_r),
\label{eq:eig_m_xxx}
\end{equation}
with \(0\le \lambda_1\le \cdots \le \lambda_r\). We then define the near-null index set
\begin{equation}
\mathcal{I}_\epsilon=
\left\{
\ell \in \{1,\ldots,r\} :
\lambda_\ell \le \epsilon\,\lambda_{\max}(M)
\right\},
\label{eq:near_null_index_xxx}
\end{equation}
where \(\epsilon=0\) recovers the exact null space when it exists. Let
\[
B_\epsilon = V_{[:,\mathcal{I}_\epsilon]}
\]
denote the basis of the corresponding exact-or-near-null subspace.

Inside this subspace, XXX retains directions that preserve positive pseudo between-class scatter. To do so, it forms
\begin{equation}
A_\epsilon = B_\epsilon^\top A B_\epsilon,
\label{eq:a_eps_xxx}
\end{equation}
and computes its eigendecomposition
\begin{equation}
A_\epsilon = R\Gamma R^\top.
\label{eq:eig_a_eps_xxx}
\end{equation}
Let \(R_+\) contain the eigenvectors associated with positive numerically nonzero eigenvalues of \(\Gamma\). The final projection matrix is
\begin{equation}
W = U_r B_\epsilon R_+ \in \mathbb{R}^{d\times L},
\label{eq:projection_xxx_revised}
\end{equation}
where \(L\) is the number of retained directions. A configuration yielding \(L=0\) is treated as invalid and is not used.

A sample \(x\in\mathbb{R}^{d}\) is projected as
\begin{equation}
z = W^\top x \in \mathbb{R}^{L}.
\label{eq:projection_sample_xxx_revised}
\end{equation}

Figure~\ref{fig:svd_vs_gs} illustrates the motivation for using SVD-based basis construction in large tabular settings.

\begin{figure}[!ht]
    \centering
    \includegraphics[width=\columnwidth]{pictures/svd_vs_gs_projection.png}
    \caption{Illustration of basis-construction stability in XXX. In large tabular settings with redundant or highly correlated samples, SVD provides a more stable basis for the reduced NFST computation than classical Gram--Schmidt orthogonalization.}
    \label{fig:svd_vs_gs}
\end{figure}

This construction coincides with classical NFST when exact null directions are available and \(\epsilon=0\), and otherwise yields a numerically stable approximate NFST projection based on near-null directions.

\subsection{Representative-Based Novelty Scoring}
\label{subsec:scoring_xxx}

After learning \(W\), each pseudo-class is represented in the projected space by its projected centroid
\begin{equation}
p_j = W^\top \mu_j,
\qquad j=1,\ldots,c.
\label{eq:representatives_xxx}
\end{equation}
Given a test sample \(x_{\text{test}}\), let
\[
z_{\text{test}} = W^\top x_{\text{test}}.
\]
Its novelty score is defined as the minimum Euclidean distance to the projected pseudo-class representatives:
\begin{equation}
s(x_{\text{test}})
=
\min_{j=1,\ldots,c}
\left\|
z_{\text{test}}-p_j
\right\|_2.
\label{eq:novelty_score_xxx_revised}
\end{equation}

This score is intentionally prototype-based. It captures multiple benign modes at low computational cost, but it does not explicitly model within-mode covariance or local density. Accordingly, XXX should be interpreted as a lightweight multi-representative detector rather than a full covariance-aware local model.

\subsection{Threshold Estimation and Decision Rule}
\label{subsec:threshold_xxx}

To convert novelty scores into binary decisions, XXX first computes training scores for benign samples:
\begin{equation}
s_i
=
\min_{j=1,\ldots,c}
\left\|
W^\top x_i - p_j
\right\|_2,
\qquad i=1,\ldots,n.
\label{eq:train_scores_xxx_revised}
\end{equation}
The decision threshold is then estimated from the empirical distribution of these scores:
\begin{equation}
\tau = Q_{1-\rho}\big(\{s_i\}_{i=1}^{n}\big),
\label{eq:tau_xxx_revised}
\end{equation}
where \(Q_{1-\rho}\) denotes the \((1-\rho)\)-quantile and \(\rho\) controls how much of the upper tail of benign training scores is ignored. The special case \(\rho=0\) reduces to the maximum training score.

The final decision rule is
\begin{equation}
f(x_{\text{test}})=
\begin{cases}
0, & \text{if } s(x_{\text{test}})\le \tau,\\
1, & \text{if } s(x_{\text{test}})> \tau.
\end{cases}
\label{eq:decision_rule_xxx_revised}
\end{equation}

Hence, XXX combines \emph{representative-based localized scoring} with a \emph{single global threshold}. Figure~\ref{fig:threshold_rewrite} illustrates this thresholding step.

\begin{figure}[!ht]
    \centering
    \includegraphics[width=\columnwidth]{pictures/threshold.png}
    \caption{Illustration of threshold estimation in XXX. After representative-based scoring is computed for benign training samples, a single global threshold is selected from the empirical score distribution using the quantile parameter \(\rho\).}
    \label{fig:threshold_rewrite}
\end{figure}

\subsection{Training and Inference Procedures}
\label{subsec:algorithms_xxx_revised}

Algorithm~\ref{alg:train_xxx_revised} summarizes the training procedure and Algorithm~\ref{alg:test_xxx_revised} summarizes inference.

\begin{algorithm}[!htbp]
\caption{Training Procedure of XXX}
\label{alg:train_xxx_revised}
\begin{algorithmic}[1]
\REQUIRE Benign training matrix \(X=[x_1,\ldots,x_n]\in\mathbb{R}^{d\times n}\), number of pseudo-classes \(c\ge 2\), near-null tolerance \(\epsilon\), threshold parameter \(\rho\)
\ENSURE Projection matrix \(W\), projected representatives \(\{p_j\}_{j=1}^{c}\), threshold \(\tau\)

\STATE Apply K-means to partition \(X\) into pseudo-classes \(\{C_1,\ldots,C_c\}\)
\STATE Compute pseudo-class centroids \(\{\mu_j\}_{j=1}^{c}\) and global mean \(\mu\)
\STATE Compute \(S_w\) and \(S_b\) using Eqs.~(\ref{eq:sw_xxx_revised})--(\ref{eq:sb_xxx_revised})
\STATE Form the centered data matrix \(X_c\) using Eq.~(\ref{eq:centered_matrix_xxx_revised})
\STATE Compute the compact SVD \(X_c = U_r \Sigma_r V_r^\top\)
\STATE Form \(M = U_r^\top S_w U_r\) and \(A = U_r^\top S_b U_r\)
\STATE Compute the eigendecomposition \(M=V\Lambda V^\top\)
\STATE Build the near-null index set \(\mathcal{I}_\epsilon\) using Eq.~(\ref{eq:near_null_index_xxx})
\STATE Form \(B_\epsilon = V_{[:,\mathcal{I}_\epsilon]}\)
\STATE Compute \(A_\epsilon = B_\epsilon^\top A B_\epsilon\) and its eigendecomposition \(A_\epsilon = R\Gamma R^\top\)
\STATE Retain the eigenvectors corresponding to positive eigenvalues of \(\Gamma\) to obtain \(R_+\)
\STATE Form the projection matrix \(W = U_r B_\epsilon R_+\)
\STATE Compute projected representatives \(p_j = W^\top \mu_j\), \(j=1,\ldots,c\)
\STATE Compute benign training scores \(\{s_i\}_{i=1}^{n}\) using Eq.~(\ref{eq:train_scores_xxx_revised})
\STATE Estimate \(\tau = Q_{1-\rho}(\{s_i\}_{i=1}^{n})\)
\RETURN \(W\), \(\{p_j\}_{j=1}^{c}\), \(\tau\)
\end{algorithmic}
\end{algorithm}

\begin{algorithm}[!htbp]
\caption{Inference Procedure of XXX}
\label{alg:test_xxx_revised}
\begin{algorithmic}[1]
\REQUIRE Test sample \(x_{\text{test}}\in\mathbb{R}^{d}\), projection matrix \(W\), projected representatives \(\{p_j\}_{j=1}^{c}\), threshold \(\tau\)
\ENSURE Predicted label

\STATE Compute the projected test representation \(z_{\text{test}} = W^\top x_{\text{test}}\)
\STATE Compute the novelty score
\[
s(x_{\text{test}})
=
\min_{j=1,\ldots,c}
\left\|
z_{\text{test}}-p_j
\right\|_2
\]
\STATE Predict \texttt{normal} if \(s(x_{\text{test}})\le \tau\); otherwise predict \texttt{anomalous}
\RETURN Predicted label
\end{algorithmic}
\end{algorithm}

\subsection{Complexity Analysis}
\label{subsec:complexity_xxx_revised}

The computational cost of XXX is divided into pseudo-class construction, projection learning, threshold estimation, and inference.

\paragraph{Pseudo-class construction.}
For \(n\) samples, \(d\) features, \(c\) pseudo-classes, and \(I\) K-means iterations, the clustering cost is
\[
O(ncdI).
\]

\paragraph{Projection learning.}
Computing centroids and dense scatter matrices requires \(O(nd + nd^2)\) time. For \(n>d\), the compact SVD of \(X_c\in\mathbb{R}^{d\times n}\) requires \(O(nd^2)\) time, and the eigendecompositions of the reduced matrices require at most \(O(d^3)\).

\paragraph{Threshold estimation.}
After \(W\) is learned, computing projected training scores requires projecting all training samples and comparing them to \(c\) representatives. This costs
\[
O(ndL + ncL).
\]

\paragraph{Overall training complexity.}
The total training complexity is therefore
\begin{equation}
O(ncdI + nd^2 + d^3 + ndL + ncL).
\label{eq:complexity_train_xxx_revised}
\end{equation}
In the common large-tabular regime \(n\gg d\), the dominant terms are typically K-means, SVD, and the computation of training scores.

\paragraph{Inference complexity.}
For a single test sample, projection requires \(O(dL)\), and comparison with \(c\) representatives requires \(O(cL)\). Hence the per-sample inference cost is
\begin{equation}
O(dL + cL).
\label{eq:complexity_test_xxx_revised}
\end{equation}

\paragraph{Memory complexity.}
During training, XXX stores the data matrix or centered data matrix, the scatter matrices, cluster assignments, the projection matrix, and the projected representatives. Under dense storage, the dominant training memory cost is
\begin{equation}
O(nd + d^2 + dL + cL).
\label{eq:complexity_memory_train_xxx}
\end{equation}
At inference time, only \(W\), \(\{p_j\}_{j=1}^{c}\), and \(\tau\) are required, giving an inference memory footprint of
\begin{equation}
O(dL + cL).
\label{eq:complexity_memory_test_xxx}
\end{equation}


\section{Experimental Evaluation}
\subsection{Experimental Setup}
\label{sec:exp_su}

We evaluate XXX through three research questions:

\begin{itemize}
    \item \textbf{RQ1:} How well does XXX perform against representative anomaly-detection baselines on public IoT datasets?
    \item \textbf{RQ2:} How robust is XXX under varying contamination levels and heterogeneous anomaly types?
    \item \textbf{RQ3:} How efficient and scalable is XXX on large tabular IoT datasets?
\end{itemize}

\subsubsection{Datasets}
\label{subsec:datasets}

For RQ1, we use four public IoT benchmarks: CICIoT2023, ToN-IoT, N-BaIoT, and BoT-IoT~\cite{neto2023ciciot2023,8490192,koroniotis2019towards,fesz-dm97-19}. For each dataset, we construct an evaluation subset of 100{,}000 instances and split it into training and test partitions with a 70:30 ratio. The training partition contains only benign samples, while the test partition contains the remaining benign samples together with attack samples. Full dataset descriptions and class summaries are provided in ~\ref{ap1:datasets}.

For RQ2, we generate four synthetic datasets derived from the real-world IoT benchmarks following the anomaly-generation protocol of~\cite{han2022adbench}, incorporating local, global, and cluster anomalies. Contamination levels of \(1\%\), \(3\%\), and \(5\%\) are used to simulate varying benign-set noise. Detailed generation rules are provided in ~\ref{ap1:ano_types}.

For RQ3, we evaluate scalability on synthetic tabular datasets with sample sizes ranging from \(10^3\) to \(10^6\) and feature dimensions ranging from \(10\) to \(10^4\), and record training time, inference time, and peak RAM usage.

\subsubsection{Baselines, Preprocessing, and Evaluation Protocol}
\label{subsec:bs_cp}

We compare XXX against 21 representative anomaly-detection baselines implemented through PyOD~\cite{zhao2019pyod}. The full list of models is provided in ~\ref{ap1:ap_baselines}. All unsupervised baselines are trained on the same benign-only training split as XXX to ensure a consistent comparison setting.

Before training and evaluation, we apply standard preprocessing steps, including duplicate removal, missing-value handling, categorical encoding, and feature normalization. For RQ1 and RQ2, the primary evaluation metrics are AUC-ROC and AUC-PR. We additionally report the F1-score in the main text, while ACC, Precision, Recall, FPR, as well as detailed thresholding rules, hyperparameter settings, and metric definitions are provided in Appendix~\ref{ap1:EM}.

\subsubsection{Implementation Details}
\label{subsec:implementation}

All experiments are conducted in Python 3.10 on Ubuntu 22.04 with an Intel i5-13400F CPU and 64GB RAM. Since no public Python implementation of the linear NFST formulation was available, we implemented XXX ourselves. Preprocessing uses Scikit-learn~\cite{pedregosa2011scikit}, and baseline methods are implemented through PyOD~\cite{zhao2019pyod}.




% Although this formulation provides strong discriminative power 
% in multi-class settings, several challenges arise when applying 
% NFST in practical one-class novelty detection problems. In 
% particular, the absence of between-class statistics may lead 
% to degeneracy of the Fisher objective, while the computation 
% of the orthonormal basis using Gram--Schmidt can suffer from 
% numerical instability when the data matrix is ill-conditioned.
% These issues motivate the need for more robust formulations 
% when applying NFST to real-world anomaly detection tasks.



% However, despite its strong discriminative property,
% NFST exhibits several limitations.

% First, the formulation focuses only on global class
% statistics represented by class means, and therefore
% does not preserve the local structure of samples
% within a class. In practice, data distributions may
% contain multiple local clusters within the same class,
% and collapsing all samples to a single point may
% destroy important structural information.

% Figure~\ref{fig:collapse_issue} illustrates the collapse
% phenomenon in the NFST projection space. When two
% classes share crossover points or overlapping affine
% hulls, NFST may map samples from multiple classes
% to the same location in the null space. This results
% in ambiguous representations and can lead to
% misclassification during testing.

% \begin{figure}[t]
% \centering
% % \includegraphics[width=\linewidth]{collapse_figure}
% \caption{Illustration of the collapsing singular-point
% problem in the NFST projection space. When classes
% contain crossover points, the NFST projection may map
% multiple classes to the same location in the null space,
% leading to ambiguous representations.}
% \label{fig:collapse_issue}
% \end{figure}

% Another limitation arises from the numerical computation
% of the projection basis. In the standard NFST algorithm,
% an orthonormal basis for the column space of the
% centered data matrix is typically obtained using the
% Gram–Schmidt orthogonalization procedure. While
% Gram–Schmidt theoretically produces an orthonormal
% basis, it is known to suffer from numerical instability
% when the input vectors are nearly linearly dependent
% or when the data matrix is ill-conditioned.

% In such cases, accumulated floating-point errors may
% produce a basis that deviates from true orthogonality,
% which in turn affects the computation of the null space
% of the matrix involved in the NFST formulation.
% Consequently, the resulting projection directions may
% be inaccurate and lead to degraded separation in the
% projected space.

% Figure~\ref{fig:gs_instability} illustrates this instability
% by comparing the embeddings obtained under different
% null-space computation strategies.

% \begin{figure}[t]
% \centering
% % \includegraphics[width=\linewidth]{instability_figure}
% \caption{Illustration of numerical instability in
% Gram–Schmidt-based null-space computation. Errors
% accumulated during orthogonalization can distort the
% resulting projection directions and reduce cluster
% separability in the embedding space.}
% \label{fig:gs_instability}
% \end{figure}

% These limitations motivate the development of improved
% NFST variants that better preserve the intrinsic data
% structure and provide more stable projections.



% The Null Foley-Sammon Transform (NFST), introduced by Guo et al.~\cite{Guo2006}, is a discriminative linear projection technique designed to maximize class separability based on the Fisher discriminant criterion. NFST maps all samples of a given class to a single point in a discriminative nullspace, achieving zero within-class variance while maintaining positive between-class variance. This section provides a detailed mathematical formulation of NFST, its key properties, and its limitations, particularly in the small sample size (SSS) case where the number of training samples \( n \) is less than the feature dimension \( d \). Consider a dataset with \( n \) samples \( \{\mathbf{x}_k\}_{k=1}^n \in \mathbb{R}^d \), belonging to \( c \) classes, where the \( i \)-th class has \( n_i \) samples with class mean \( \mathbf{m}_i \), and \( \mathbf{m} \) is the overall mean of the dataset. For a sample \( \mathbf{x}_k \) in the \( j \)-th class, define the centered sample as \( \overline{\mathbf{x}}_k = \mathbf{x}_k - \mathbf{m}_j \). The within-class scatter matrix \( \mathbf{S}_w \in \mathbb{R}^{d \times d} \) and between-class scatter matrix \( \mathbf{S}_b \in \mathbb{R}^{d \times d} \) are defined as:
% \[
% \mathbf{S}_w = \sum_{k=1}^n \overline{\mathbf{x}}_k \overline{\mathbf{x}}_k^\top, \quad \mathbf{S}_b = \sum_{i=1}^c n_i (\mathbf{m}_i - \mathbf{m})(\mathbf{m}_i - \mathbf{m})^\top.
% \]
% The total scatter matrix is given by \( \mathbf{S}_t = \mathbf{S}_w + \mathbf{S}_b \).

% The Foley-Sammon Transform (FST) ~\cite{Foley1975} seeks an optimal set of orthonormal discriminant vectors \( \mathbf{w}_i \in \mathbb{R}^d \) that maximize the Fisher criterion:
% % \[
% % J(\mathbf{w}_i) = \frac{\mathbf{w}_i^\top \mathbf{S}_b \mathbf{w}_i}{\mathbf{w}_i^\top \mathbf{S}_w \mathbf{w}_i}, \quad \text{s.t.} \quad \mathbf{w}_i^\top \mathbf{w}_j = \delta_{ij}, \quad j = 1, \ldots, i-1,
% % \]

% \begin{equation}
% \begin{aligned}
% J(\mathbf{w}_i) &= \frac{\mathbf{w}_i^\top \mathbf{S}_b \mathbf{w}_i}{\mathbf{w}_i^\top \mathbf{S}_w \mathbf{w}_i} \\
% \text{subject to} \quad & \mathbf{w}_i^\top \mathbf{w}_j = \delta_{ij}, \quad j = 1, \ldots, i-1
% \end{aligned}
% \label{eq:FisherCri}
% \end{equation}


% where \( \delta_{ij} \) is the Kronecker delta. The NFST extends FST by utilizing the nullspace of \( \mathbf{S}_w \). Specifically, NFST seeks discriminant vectors \( \mathbf{w} \), termed Null Projecting Directions (NPDs), where the numerator and denominator in Eq.\eqref{eq:FisherCri} satisfy: 
% \[
% \mathbf{w}^\top \mathbf{S}_w \mathbf{w} = 0, \quad \mathbf{w}^\top \mathbf{S}_b \mathbf{w} > 0.
% \]
% Since \( \mathbf{S}_w \) is positive semi-definite, the condition \( \mathbf{w}^\top \mathbf{S}_w \mathbf{w} = 0 \) implies \( \mathbf{S}_w \mathbf{w} = \mathbf{0} \) (by Lemma 1 in~\cite{Ali2023}), meaning \( \mathbf{w} \) lies in the nullspace of \( \mathbf{S}_w \), denoted:
% \[
% \mathbf{N}_w = \{\mathbf{v} \in \mathbb{R}^d \mid \mathbf{S}_w \mathbf{v} = \mathbf{0}\}.
% \]
% The condition \( \mathbf{w}^\top \mathbf{S}_b \mathbf{w} > 0 \) ensures discriminative capability, equivalent to \( \mathbf{w}^\top \mathbf{S}_t \mathbf{w} > 0 \), as \( \mathbf{S}_t = \mathbf{S}_w + \mathbf{S}_b \). Thus, NPDs lie in the intersection of \( \mathbf{N}_w \) and the orthogonal complement of the nullspace of \( \mathbf{S}_t \), i.e., \( \mathbf{N}_t^\perp \), where:
% \[
% \mathbf{N}_t = \{\mathbf{v} \in \mathbb{R}^d \mid \mathbf{S}_t \mathbf{v} = \mathbf{0}\}.
% \]
% The projection matrix \( \mathbf{W} = [\mathbf{w}_1, \mathbf{w}_2, \ldots, \mathbf{w}_L] \in \mathbb{R}^{d \times L} \) is formed by an orthonormal basis of subspace \( \mathbf{N}_t^\perp \cap \mathbf{N}_w \), and the NFST is defined as:
% \[
% \mathbf{y}_k = \mathbf{W}^\top \mathbf{x}_k.
% \]
% This transformation ensures that all samples of a class collapse to a single point in the discriminative nullspace, as illustrated in Fig.~\ref{fig:cluster_impact} (adapted from~\cite{Ali2023}).



% To compute the NPDs, let \( \mathbf{X}_t \in \mathbb{R}^{d \times n} \) contain the zero-mean centered data \( \{\mathbf{x}_k - \mathbf{m}\}_{k=1}^n \). The rowspace of \( \mathbf{S}_t = \mathbf{X}_t \mathbf{X}_t^\top \) is the column space of \( \mathbf{X}_t \). Assuming \( \mathbf{X}_t \) has rank \( r \), Gram-Schmidt orthonormalization on \( \{\mathbf{x}_k - \mathbf{m}\}_{k=1}^n \) yields an orthonormal basis \( \mathbf{Q} = [\mathbf{q}_1, \ldots, \mathbf{q}_r] \in \mathbb{R}^{d \times r} \) for \( \mathbf{N}_t^\perp \). A vector \( \mathbf{w} \in \mathbf{N}_t^\perp \) can be expressed as:
% \[
% \mathbf{w} = \mathbf{Q} \boldsymbol{\beta}, \quad \boldsymbol{\beta} \in \mathbb{R}^r.
% \]
% For \( \mathbf{w} \in \mathbf{N}_w \), we require \( \mathbf{S}_w \mathbf{w} = \mathbf{S}_w \mathbf{Q} \boldsymbol{\beta} = \mathbf{0} \), which implies:
% \[
% \boldsymbol{\beta}^\top \mathbf{Q}^\top \mathbf{S}_w \mathbf{Q} \boldsymbol{\beta} = 0 \quad \Rightarrow \quad \mathbf{Q}^\top \mathbf{S}_w \mathbf{Q} \boldsymbol{\beta} = \mathbf{0}.
% \] 
% Using Lemma 1 in~\cite{Ali2023}. The nullspace of \( \mathbf{Q}^\top \mathbf{S}_w \mathbf{Q} \), with dimension \( L \), provides orthonormal vectors \( \boldsymbol{\beta}_1, \ldots, \boldsymbol{\beta}_L \). The NPDs are then \( \mathbf{w}_i = \mathbf{Q} \boldsymbol{\beta}_i \), and the projection matrix is:
% \[
% \mathbf{W} = \mathbf{Q} \mathbf{B}, \quad \mathbf{B} = [\boldsymbol{\beta}_1, \ldots, \boldsymbol{\beta}_L] \in \mathbb{R}^{r \times L}.
% \]
% The standard algorithm (s-NFST) for computing NFST is summarized in Algorithm ~\ref{alg:s-nfst} ~\cite{Ali2023}.

% \begin{algorithm}
% \caption{s-NFST Algorithm}
% \label{alg:s-nfst}
% \begin{algorithmic}
% \REQUIRE Data \( \mathbf{X} = [\mathbf{x}_1, \ldots, \mathbf{x}_n] \in \mathbb{R}^{d \times n} \) with labels
% \ENSURE Projection matrix \( \mathbf{W} \in \mathbb{R}^{d \times L} \), mapped data \( \mathbf{Y} = [\mathbf{y}_1, \ldots, \mathbf{y}_n] \in \mathbb{R}^{L \times n} \)
% \STATE Compute \( \mathbf{S}_w \) using \( \mathbf{X} \) and labels
% \STATE \( \mathbf{Q} \gets \) Gram-Schmidt orthonormalization on \( \{\mathbf{x}_k - \mathbf{m}\}_{k=1}^n \)
% \STATE \( \mathbf{B} \gets \) Nullspace basis of \( \mathbf{Q}^\top \mathbf{S}_w \mathbf{Q} \)
% \STATE \( \mathbf{W} = \mathbf{Q} \mathbf{B} \), \( \mathbf{Y} = \mathbf{W}^\top \mathbf{X} \)
% \end{algorithmic}
% \end{algorithm}




% In the s-NFST algorithm, \( \mathbf{Q} \in \mathbb{R}^{d \times r} \) is an orthonormal basis for the column space of the zero-mean centered data matrix \( \mathbf{X}_t \in \mathbb{R}^{d \times n} \), where \( r = \operatorname{rank}(\mathbf{X}_t) \leq \min(n, d) = d \) (typically \( r = d \) in the LSS case under Assumption 1(b) ~\cite{Ali2023}). Gram-Schmidt orthonormalizes the columns of \( \mathbf{X}_t \), but numerical instability can produce an ill-defined \( \mathbf{Q} \) (e.g., non-orthogonal or near-zero columns), disrupting the nullspace computation of \( \mathbf{Q}^\top \mathbf{S}_w \mathbf{Q} \). 

% For \( \mathbf{Q} \) to be ill-defined (interpreted as numerically invalid, akin to ``\( \mathbf{Q} = \mathbf{0} \)'') while \( \mathbf{U}_r \neq \mathbf{0} \), \( \mathbf{X}_t \neq \mathbf{0} \) (i.e., \( r \geq 1 \)) is required, as \( \mathbf{X}_t = \mathbf{0} \) would imply both \( \mathbf{Q} = \mathbf{0} \) and \( \mathbf{U}_r = \mathbf{0} \) (no column space). The following scenarios in the LSS case (\( n > d \)) can cause Gram-Schmidt to fail:

% \begin{definition}[Highly Collinear Data]
% If the columns of \( \mathbf{X}_t \) are nearly linearly dependent (high condition number, i.e., ratio of largest to smallest non-zero singular value is large), Gram-Schmidt accumulates numerical errors, producing a \( \mathbf{Q} \) with non-orthogonal or near-zero columns. SVD computes \( \mathbf{U}_r \) accurately, as it explicitly handles small singular values, ensuring a valid basis.
% \end{definition}

% \begin{definition}[Ill-Conditioned \( \mathbf{X}_t \)]
% If \( \mathbf{X}_t \) has very small non-zero singular values (e.g., close to machine precision), Gram-Schmidt may treat them as zero due to finite precision, leading to an incorrect or collapsed basis. SVD identifies all non-zero singular values, ensuring \( \mathbf{U}_r \) spans the correct subspace.
% \end{definition}

% \begin{definition}[Large \( n \) with Numerical Errors]
% For large \( n \), Gram-Schmidt’s sequential orthonormalization of \( n \) columns amplifies floating-point errors, especially if \( \mathbf{X}_t \) is poorly scaled. SVD, implemented in robust libraries (e.g., LAPACK), maintains orthogonality, providing a stable \( \mathbf{U}_r \).
% \end{definition}



% Draft : Standard global projection methods fail here because they attempt to collapse this complex mixture into a single point, inevitably destroying the local covariance and mean-shift information required to distinguish these specific anomaly types.


% \begin{itemize}
%     \item \textbf{Loss of local structure information}: 
%     In OCND, all training samples belong to a single normal class. NFST therefore collapses the entire training set into a single point in the projected space, implicitly assuming a unimodal normal distribution. This assumption disregards the multi-modal structure commonly observed in IoT traffic and increases sensitivity to noise and contamination~\cite{ali2022theoretical, nguyen2025nnfst}.
    
%     \item \textbf{Computational complexity}: 
%     Standard NFST relies on eigen-decomposition and matrix operations with computational complexity up to $\mathcal{O}(N^3)$, limiting scalability in trainning process for large IoT datasets~\cite{ali2022theoretical}.
% \end{itemize}







%     % Moreover, even after clustering the training data into multiple groups, classical NFST still computes the distance between a projected test point and a single global singular point, and compares it against a fixed threshold—rather than learning local thresholds adapted to each cluster. 
    
    
%     % (2) when normal data exhibit multiple distinct sub-clusters, collapsing them into a single projection point leads to poor representation and an increased false negative rate, as the model cannot distinguish between intra-class variations and true anomalies.
    
% % \subsection{Proposed Solution: XXX}

% % Acknowledging the aforementioned limitations of traditional NFST and the research gap in applying subspace projection methods to one-class novelty detection (OCND) in IoT-based intrusion detection systems (IDSs), we propose a novel and lightweight algorithm named \textbf{Local One-Class Null Foley-Sammon Transformation (XXX)}. This approach adapts and enhances the original NFST framework to suit single-class, and noisy IoT traffic data.

% % The key improvements introduced in XXX are as follows:

% % \begin{itemize}
% %     \item To address the first challenge: 
% %         \begin{itemize}
% %             \item (1) Firstly, our proposed XXX framework introduces a K-means clustering step on the normal training data to generate pseudo-classes. Traditional NFST is inherently a multi-class method and lacks compatibility with single-class scenarios, making it unsuitable for many real-world IoT settings where only normal (benign) traffic is labeled during training. To overcome this, XXX introduces a clustering step using K-Means to partition the normal da`ta into pseudo-classes. This transformation enables the formulation of intra-class and inter-class scatter matrices—critical to NFST—while remaining within the one-class learning paradigm.
% %             \item (2) Secondly,  XXX adopts a localized decision strategy by computing anomaly scores based on cluster-specific proximity instead of using a single global threshold for all test samples, which fails to capture the local structure of clustered training data. Specifically, we define an anomaly score for a test sample \( x_{\text{test}} \). By utilizing cluster-aware scoring rather than global thresholding, XXX mitigates the overgeneralization inherent in classical NFST and enhances its sensitivity to fine-grained local structures in the data.            
% %             % ~\cite{nguyen2025nnfst, hu2025efficient}.
% %         \end{itemize}
        

% %     \item In response to the second challenge,  This modification effectively transforms the one-class training setting into a pseudo multi-class scenario, allowing the model to construct meaningful within-class and between-class scatter matrices even when the number of features is limited. Notably, the number of pseudo-classes becomes a tunable hyperparameter, which we observe to have a direct influence on the generation of Null Projecting Directions (NPDs)—the discriminative directions used to map original data points into the null space. By incorporating K-means clustering into the model pipeline, we ensure that valid NPDs can always be computed, thereby eliminating the small-sample-size constraint of traditional NFST. The theoretical justification for the guaranteed existence of NPDs under this formulation is provided in Section~\ref{sec:method}.

% %     \item For the final challenge, our proposed XXX replaces the Gram-Schmidt orthonormalization process used in classical NFST with \textbf{Singular Value Decomposition (SVD)}. While Gram-Schmidt is a standard method for generating orthonormal bases, it often suffers from numerical instability and accumulated rounding errors, especially in high-dimensional or nearly linearly dependent feature spaces—conditions common in IoT traffic data. In contrast, SVD offers a stable and robust approach to subspace projection by decomposing matrices in a way that is less sensitive to floating-point inaccuracies~\cite{golub2013matrix}. This modification enhances both the computational robustness and efficiency of XXX, enabling it to scale effectively to large feature spaces without sacrificing precision or reliability.

% % \end{itemize}

% % To the best of our knowledge, this is the first study to formulate a null-space-based OCND method for intrusion detection in IoT networks. Our contributions lie in both the methodological innovation and the practical demonstration of XXX’s robustness across various types of anomalies—including global, cluster, and dependency outliers—in multiple benchmark IoT datasets. This positions XXX as a viable lightweight IDS solution for real-world deployments in edge and fog computing environments.


% In this work, we focus on the OCND setting where previously unseen attacks are assumed to generate observable structural deviations in network traffic features. 
% Our goal is not to detect anomalies that lack measurable signals, but to design a lightweight method that exploits such signals when they are present, while remaining robust to annotation noise and suitable for deployment on edge devices.

% Existing OCND approaches for IDS can be broadly categorized into three families: hypersphere-based distance estimation, subspace projection methods, and representation learning approaches. 
% Hypersphere-based methods, including One-Class SVM (OCSVM), DADVDD, and Deep SVDD~\cite{yang2022efficient, fu2024dense, ruff2018deep}, learn compact decision boundaries around normal data but suffer from poor scalability, with computational complexity reaching $\mathcal{O}(N^3)$~\cite{KANG2019772}, which limits their applicability in large-scale IoT deployments. 
% Subspace projection techniques aim to enhance separability through dimensionality reduction, but often rely on assumptions such as $n < d$ that are violated in high-volume IoT traffic scenarios. 
% Representation learning methods, including autoencoders, variational autoencoders, and GAN-based models~\cite{perera2019ocgan, LUNAR}, detect anomalies through reconstruction errors or learned decision boundaries; however, when trained on imperfect data, these models may reconstruct anomalous or noisy samples effectively, thereby weakening the anomaly signal while imposing substantial computational costs~\cite{beggel2019robust, jodelka2021adaptive}.

% Consequently, existing approaches face inherent trade-offs and fail to simultaneously satisfy the core requirements of OCND in IoT environments: robustness under annotation noise, sensitivity to diverse anomaly structures, scalability to high-dimensional traffic data, and feasibility on resource-constrained devices. 
% To address these limitations under explicit observability assumptions, we propose the \textbf{Local One-Class Null Foley--Sammon Transformation (XXX)}, Tận dụng cơ chế nén để giải quyết noise, a lightweight method designed to exploit locally and globally observable anomaly structures in IoT intrusion detection systems.

% \subsection{Null Foley--Sammon Transformation}

% The Null Foley--Sammon Transformation (NFST)~\cite{Guo2006} is a discriminant analysis technique originally developed for Multi-class Novelty Detection (MND). 
% NFST maximizes the Fisher criterion by projecting data into a null space in which intra-class variance collapses to zero while preserving non-zero between-class variance, resulting in theoretically infinite class separability. 
% This property is particularly effective in domains such as computer vision, where class labels are abundant and class boundaries are well-defined~\cite{ayad2024efficient, perera2021one}.

% However, directly applying NFST to One-Class Novelty Detection in IoT settings is non-trivial and introduces several challenges that limit the validity of anomaly detection claims:

% \begin{itemize}
%     \item \textbf{Loss of Local Structure Information}: 
%     In OCND, all training samples belong to a single normal class. 
%     NFST therefore collapses the entire training set into a single point in the projected space, implicitly assuming a unimodal normal distribution. 
%     This assumption neglects the multi-modal structure commonly observed in IoT traffic and amplifies sensitivity to noise and contaminated samples~\cite{ali2022theoretical, nguyen2025nnfst}.
    
%     \item \textbf{Small Sample Size (SSS) Constraint}: 
%     NFST requires the number of samples to be smaller than the feature dimensionality ($n < d$) to ensure a valid null space projection~\cite{ali2022theoretical}. 
%     This condition is rarely satisfied in IoT datasets, where traffic logs typically contain far more samples than features ($n \gg d$), leading to numerical instability and invalid projections in large-scale deployments~\cite{nguyen2025nnfst}.
    
%     \item \textbf{Computational Complexity}: 
%     Standard NFST relies on expensive matrix decompositions with computational complexity up to $\mathcal{O}(N^3)$, rendering it unsuitable for large IoT datasets and real-time intrusion detection on edge devices~\cite{ali2022theoretical}.
% \end{itemize}



%     % Moreover, even after clustering the training data into multiple groups, classical NFST still computes the distance between a projected test point and a single global singular point, and compares it against a fixed threshold—rather than learning local thresholds adapted to each cluster. 
    
    
%     % (2) when normal data exhibit multiple distinct sub-clusters, collapsing them into a single projection point leads to poor representation and an increased false negative rate, as the model cannot distinguish between intra-class variations and true anomalies.
    
% % \subsection{Proposed Solution: XXX}

% % Acknowledging the aforementioned limitations of traditional NFST and the research gap in applying subspace projection methods to one-class novelty detection (OCND) in IoT-based intrusion detection systems (IDSs), we propose a novel and lightweight algorithm named \textbf{Local One-Class Null Foley-Sammon Transformation (XXX)}. This approach adapts and enhances the original NFST framework to suit single-class, and noisy IoT traffic data.

% % The key improvements introduced in XXX are as follows:

% % \begin{itemize}
% %     \item To address the first challenge: 
% %         \begin{itemize}
% %             \item (1) Firstly, our proposed XXX framework introduces a K-means clustering step on the normal training data to generate pseudo-classes. Traditional NFST is inherently a multi-class method and lacks compatibility with single-class scenarios, making it unsuitable for many real-world IoT settings where only normal (benign) traffic is labeled during training. To overcome this, XXX introduces a clustering step using K-Means to partition the normal da`ta into pseudo-classes. This transformation enables the formulation of intra-class and inter-class scatter matrices—critical to NFST—while remaining within the one-class learning paradigm.
% %             \item (2) Secondly,  XXX adopts a localized decision strategy by computing anomaly scores based on cluster-specific proximity instead of using a single global threshold for all test samples, which fails to capture the local structure of clustered training data. Specifically, we define an anomaly score for a test sample \( x_{\text{test}} \). By utilizing cluster-aware scoring rather than global thresholding, XXX mitigates the overgeneralization inherent in classical NFST and enhances its sensitivity to fine-grained local structures in the data.            
% %             % ~\cite{nguyen2025nnfst, hu2025efficient}.
% %         \end{itemize}
        

% %     \item In response to the second challenge,  This modification effectively transforms the one-class training setting into a pseudo multi-class scenario, allowing the model to construct meaningful within-class and between-class scatter matrices even when the number of features is limited. Notably, the number of pseudo-classes becomes a tunable hyperparameter, which we observe to have a direct influence on the generation of Null Projecting Directions (NPDs)—the discriminative directions used to map original data points into the null space. By incorporating K-means clustering into the model pipeline, we ensure that valid NPDs can always be computed, thereby eliminating the small-sample-size constraint of traditional NFST. The theoretical justification for the guaranteed existence of NPDs under this formulation is provided in Section~\ref{sec:method}.

% %     \item For the final challenge, our proposed XXX replaces the Gram-Schmidt orthonormalization process used in classical NFST with \textbf{Singular Value Decomposition (SVD)}. While Gram-Schmidt is a standard method for generating orthonormal bases, it often suffers from numerical instability and accumulated rounding errors, especially in high-dimensional or nearly linearly dependent feature spaces—conditions common in IoT traffic data. In contrast, SVD offers a stable and robust approach to subspace projection by decomposing matrices in a way that is less sensitive to floating-point inaccuracies~\cite{golub2013matrix}. This modification enhances both the computational robustness and efficiency of XXX, enabling it to scale effectively to large feature spaces without sacrificing precision or reliability.

% % \end{itemize}

% % To the best of our knowledge, this is the first study to formulate a null-space-based OCND method for intrusion detection in IoT networks. Our contributions lie in both the methodological innovation and the practical demonstration of XXX’s robustness across various types of anomalies—including global, cluster, and dependency outliers—in multiple benchmark IoT datasets. This positions XXX as a viable lightweight IDS solution for real-world deployments in edge and fog computing environments.


% \subsection{Main Contributions}
% To overcome these challenges, we propose \textbf{Local One-Class Null Foley-Sammon Transformation}, a pioneering single-model algorithm for OCND in IoT networks. XXX adapts NFST for one-class data using distance-based clustering to generate pseudo-classes and Singular Value Decomposition (SVD) for stable subspace projection, excelling in large-sample scenarios ($n \gg d$). The main contributions of this paper are summarized as follows:: 
    
% \begin{itemize}
%     % Contribution 1: FOCUS ON TRAINING PHASE - Technical Novelty
%     % Nhấn mạnh: SVD thay thế Gram-Schmidt để xử lý Large-Scale Data (n >> d)
%     % draft : integrates \textbf{distance-based clustering} to generate pseudo-classes
%     \item We propose \textbf{XXX}, a novel one-class anomaly detection framework tailored for IoT traffic. Unlike traditional NFST, which is constrained by small sample sizes, our approach partitions normal traffic into K pseudo-classes to enable null-space discrimination and employs \textbf{Singular Value Decomposition (SVD)} to replace the numerically unstable Gram-Schmidt process. This formulation specifically addresses the challenge of large-scale data ($n \gg d$) in IoT, ensuring the extraction of Null Projecting Directions (NPDs) while preserving global and local data structures in a single-class setting.

%     % Contribution 2: FOCUS ON INFERENCE PHASE - Methodological Novelty
%     % Nhấn mạnh: Robustness (chống nhiễu) + Phân biệt Scalable Training vs Lightweight Inference
%     \item We introduce a \textbf{cluster-aware localized scoring strategy} that calculates anomaly scores based on the minimum distance to cluster-specific singular points rather than a global center. This mechanism significantly enhances robustness against \textbf{annotation errors} and diverse anomaly types (including local, cluster, and global outliers) by capturing multi-modal distributions. Furthermore, by optimizing the subspace projection via SVD, the algorithm achieves a \textbf{scalable training complexity of $\mathcal{O}(nd^2)$} and enables \textbf{lightweight inference}, making it highly suitable for real-time deployment on resource-constrained devices.

%     % Contribution 3: Validation
%     % Nhấn mạnh: Kết quả SOTA trên 4 tập dữ liệu
%     \item \textbf{Empirical validation}: Extensive experiments on four benchmark IoT datasets—\textbf{CICIoT2023}, \textbf{N-BaIoT}, \textbf{ToNIoT}, and \textbf{BoTIoT}—demonstrate XXX’s state-of-the-art performance. The method achieves AUC-ROC scores of up to 99.61\%, outperforming 18 baseline methods in detecting complex cyberattacks~\cite{neto2023ciciot2023, 8490192, koroniotis2019towards, fesz-dm97-19}.
% \end{itemize}
% ========================



% Draft : Consequently, complex deep learning models (requiring heavy GPU matrix operations) or iterative optimization methods (undefined runtime) are unsuitable. The solution requires a deterministic, closed-form inference mechanism with complexity strictly bounded by $\mathcal{O}(d^2)$ or $\mathcal{O}(dL)$ (where $L \ll d$), executable on standard edge CPUs (e.g., ARM Cortex).
% \textbf{3. High-Throughput and Low-Latency Constraints (Real-time Requirement):}
% Unlike static data analysis, IoT-IDS operates in a streaming environment where network throughput strictly dictates the processing deadline. In the context of 5G-Advanced and early 6G scenarios~\cite{deshmukh20256g, zhao2024ngmn}, peak data rates at \textbf{IoT aggregation points (Gateways)} or high-bandwidth applications (e.g., XR, holographic comms) are projected to reach $1 \text{--} 10$ Gbit/s. This imposes a severe \textbf{arrival rate constraint}.
% \textbf{3. High-Throughput and Low-Latency Constraints (Real-time Requirement):}
% Unlike static data analysis, IoT-IDS operates in a streaming environment where network throughput determines the processing deadline. According to recent vision papers for 5G-Advanced (Rel-18/19) and early 6G scenarios~\cite{deshmukh20256g, zhao2024ngmn}, user-experienced data rates are projected to reach $1 \text{--} 10$ Gbit/s, scaling up to $100$ Gbit/s in the 2028–2030 timeframe. This imposes a strict \textbf{arrival rate constraint}.

% Let $\lambda$ be the packet arrival rate (packets/second). To prevent buffer overflow and ensure zero-latency capability during attacks like DDoS, the inference time $\mathcal{T}_{inf}$ for each sample $x$ must satisfy:
% \begin{equation}
%     \mathcal{T}_{inf}(f(x)) \leq \frac{1}{\lambda}
%     \label{eq:latency_constraint}
% \end{equation}
% At 10 Gbit/s with an average packet size of 500 bytes, the inter-arrival time is roughly $0.4 \mu s$. This necessitates an ultra-lightweight decision function $f(x)$. 
% Let $\lambda$ be the packet arrival rate (packets/second). To prevent buffer overflow and ensure wire-speed capability during volumetric attacks (e.g., DDoS), the inference time $\mathcal{T}_{inf}$ for each sample $x$ is bounded by:
% \begin{equation}
%     \mathcal{T}_{inf}(f(x)) \leq \frac{1}{\lambda} = \frac{\mathcal{L}_{pkt}}{\mathcal{R}_{line}}
%     \label{eq:latency_constraint}
% \end{equation}
% where $\mathcal{L}_{pkt}$ is the average packet size and $\mathcal{R}_{line}$ is the line rate.
% Consider a high-throughput scenario of 10 Gbit/s with a conservative average packet size of 500 bytes (based on IMIX traffic profiles). The inter-arrival time is approximately:
% \[
% \mathcal{T}_{inf} \leq \frac{500 \times 8 \text{ bits}}{10 \times 10^9 \text{ bits/s}} = 0.4 \mu s
% \]
% Even stricter constraints apply for smaller packets (e.g., 64 bytes $\rightarrow$ $0.05 \mu s$). This physical limit renders complex deep learning models (typically requiring milliseconds on CPUs) or iterative optimization methods unsuitable for line-rate detection. The solution necessitates a deterministic, closed-form inference mechanism with complexity strictly bounded by $\mathcal{O}(d^2)$ or $\mathcal{O}(dL)$ (where $L \ll d$), capable of execution on standard edge processors (e.g., ARM Cortex) within a few clock cycles.

% == ======  =      =              =  



% This paper addresses the challenge of anomaly and noise detection in IoT environments, where data is typically high-dimensional, non-stationary, and vulnerable to various irregularities \cite{al2021review}. Given a set of $N$ normal training samples, $x^{(\mathrm{train})}_1, \ldots, x^{(\mathrm{train})}_N \in \mathbb{R}^d$, and $M$ testing samples, $x^{(\mathrm{test})}_1, \ldots, x^{(\mathrm{test})}_M \in \mathbb{R}^d$, each testing sample may belong to one of two categories: normal, anomalous and noisy. The goal is to assign an anomaly score to each testing instance, where a low score indicates normal behavior and a high score signifies either an anomaly or a noisy observation.
% Problem Statement: How can we accurately detect various types of anomalies and effectively handle noisy data in real-world IoT applications, by assigning a reliable anomaly score to each test sample based on its distance from normal data points?
% Example : Outlier Detection in IoT Networks} In IoT systems, outliers can take different forms depending on the nature of the data and the attack scenario. For example, consider a smart home network where multiple devices are compromised and exhibit abnormal communication patterns. These devices may form a dense region of abnormal behavior that is isolated from the rest of the network, representing a "cluster anomaly". On the other hand, a single sensor in an industrial IoT system may report a sudden spike in temperature, deviating significantly from the normal operating range, but not necessarily belonging to a broader attack. This would be an example of a "local anomaly", detectable only with respect to the neighboring sensor data. Existing anomaly detection algorithms, such as One-Class SVM (OCSVM) and Local Outlier Factor (LOF), often specialize in detecting one type of anomaly but struggle to generalize across different types of outliers\cite{goldstein2016comparative}. A robust detection system must be capable of handling both global and local anomalies simultaneously.


% =======================================================================
% =======================================================================






% ========================


\newpage

\subsection{Results and Analysis}
This section focuses on evaluating the efficacy of XXX, presenting a comprehensive analysis of its performance across various experimental setups to validate its effectiveness in anomaly detection within IoT environments. The evaluation leverages multiple datasets and advanced analytical techniques to provide robust insights into the method's capabilities.


\label{subsec:exp_rs}
\subsubsection{Detection Performance on IoT Benchmarks (RQ1)}
\label{ssubsec:RQ1}
This section comprehensively evaluates the performance of the proposed XXX method across four distinct IoT datasets—BoIoT, CICIoT2023, N-BaIoT, and ToN-IoT—using multiple metrics, as detailed in Table~\ref{tab:iot_metrics}. The evaluation relies on Area Under the Receiver Operating Characteristic Curve (AUC-ROC), Area Under the Precision-Recall Curve (AUC-PR), Matthews Correlation Coefficient (MCC), Accuracy, F1 Score, Precision, and Recall, highlighting the method's robustness. AUC-ROC scores reach 99.39\% for N-BaIoT and 97.10\% for BoIoT, indicating strong classification performance, while AUC-PR exceeds 98\%, with CICIoT2023 achieving 99.84\%, reflecting excellent precision-recall balance. MCC peaks at 0.9865 for N-BaIoT, and Accuracy hits 99.77\% for the same dataset, demonstrating high reliability. F1 scores like 0.9987 for N-BaIoT and 0.9692 for BoIoT, alongside Precision of 0.9998 for CICIoT2023 and Recall of 0.9998 for N-BaIoT, underscore a balanced and effective anomaly detection capability across diverse IoT environments.

\begin{table*}[!ht]
\caption{Performance of the method on four IoT datasets using AUC-ROC, AUC-PR, MCC, Accuracy, F1, Precision, and Recall.}
\centering
\resizebox{\textwidth}{!}{%
\begin{tabular}{|c|c|c|c|c|c|c|c|}
\hline
\textbf{Dataset} & \textbf{AUC-ROC (\%)} & \textbf{AUC-PR (\%)} & \textbf{MCC} & \textbf{ACC (\%)} & \textbf{F1} & \textbf{Precision} & \textbf{Recall} \\
\hline
BoIoT      & 97.10 & 99.52 & 0.7378 & 94.53 & 0.9692 & 0.9891 & 0.9502 \\
CICIoT2023 & 94.93 & 99.84 & 0.3845 & 85.82 & 0.9211 & 0.9998 & 0.8540 \\
N-BaIoT    & 99.39 & 99.79 & 0.9865 & 99.77 & 0.9987 & 0.9977 & 0.9998 \\
ToN-IoT    & 94.63 & 98.74 & 0.7349 & 93.76 & 0.9641 & 0.9809 & 0.9478 \\
\hline
\end{tabular}
}
\label{tab:iot_metrics}
\end{table*}


% Required packages in the preamble:
% \usepackage[table]{xcolor}
% \usepackage{multirow}
% \usepackage{adjustbox}




Further, Critical Difference (CD) diagrams in Figure~\ref{fig:cd_map} compare XXX with baselines across the datasets, evaluating AUC-PR and AUC-ROC without relying on static tests that have been removed. The diagrams present average scores and ranks computed over the four datasets, with XXX leading in AUC-PR at 99.47\% (rank 1.00), far surpassing ECOD at 83.93\% (rank 16.75) and DeepSVDD at 77.83\% (rank 9.00), while KNN at 99.03\% (rank 3.12) and LOF at 98.74\% (rank 4.50) follow closely. In AUC-ROC, XXX achieves 99.61\% (rank 1.00), outpacing ECOD at 89.16\% (rank 16.75) and DeepSVDD at 77.83\% (rank 9.00), with KNN at 99.13\% (rank 3.25) and LOF at 98.74\% (rank 4.25) trailing, emphasizing its superior detection accuracy due to robust feature selection and transformation.





\begin{figure}[!ht]
\centering
\begin{subfigure}{\columnwidth}
\includegraphics[width=\columnwidth]{pictures/CDDiagram_AUCPR.png}
\caption{AUCPR}
\label{fig:cd_aucpr}
\end{subfigure}%
\hfill
\begin{subfigure}{\columnwidth}
\includegraphics[width=\columnwidth]{pictures/CDDiagram_AUCROC.png}
\caption{AUCROC}
\label{fig:cd_aucroc}
\end{subfigure}
\caption{AUC-ROC and AUC-PR performance comparison across four IoT datasets using Critical Difference diagrams}
\label{fig:cd_map}
\end{figure}


\subsubsection{Robustness Against Contamination and Anomaly Types (RQ2)}
\label{ssubsec:RQ2}

This section evaluates the detection performance of the proposed XXX algorithm across various anomaly types and noise levels, comparing it against several baseline methods. The performance is quantified using the Area Under the Receiver Operating Characteristic Curve (AUC-ROC), a widely accepted metric for assessing the ability of a model to distinguish between normal and anomalous instances. The results are visualized in Figure~\ref{fig:noise_vs_tp}, which comprises two subfigures: Figure~\ref{fig:noise_levels} for noise level analysis and Figure~\ref{fig:type_levels} for anomaly type analysis.

Noise Level and Anomaly Type Analysis: Figure~\ref{fig:noise_levels} illustrates the average AUC-ROC performance under different noise levels (1\%, 3\%, and 5\%), representing the proportion of corrupted data that tests the algorithms' robustness. XXX (AVG-Rank: 1.00) leads with scores around 90\% at 1\% noise, dropping slightly to 85\% at 5\% noise, showcasing exceptional stability. KNN (AVG-Rank: 2.83) follows with a strong 88\% at 1\% noise, declining to 83\% at 5\% noise, indicating good adaptability with a noticeable drop. In contrast, DeepSVDD (AVG-Rank: 8.17) lags with scores around 75\% across noise levels, reflecting lower robustness.  \\
Figure~\ref{fig:type_levels} compares performance across Global Type, Cluster Type, and Local Type anomaly patterns. XXX (AVG-Rank: 1.50) excels with near 100\% for Global and Cluster Types and 92.5\% for Local Type, demonstrating versatility. LOF (AVG-Rank: 6.58) achieves around 100\% for Global and Cluster Types but drops to 90\% for Local Type, while KNN (AVG-Rank: 7.58) reaches 100\% for Global Type but falls to 82.5\% for Local Type, showing weakness in localized anomalies.\\
Detailed Insights: The bar plots collectively demonstrate that XXX outperforms the baseline methods across both noise levels and anomaly types, as reflected in its lowest average ranks (1.00 for noise levels and 1.50 for anomaly types). This superiority is attributed to its robust feature selection and transformation process, which effectively mitigates the impact of noise and adapts to various anomaly patterns. 

\begin{figure*}[!ht]
\centering
\begin{subfigure}{0.48\textwidth}    
    \includegraphics[width=\linewidth]{pictures/noise_levels.png} 
    \caption{Noise Levels: Average AUC-ROC performance across different noise percentages for XXX and top baseline methods.}
    \label{fig:noise_levels}
\end{subfigure}\hfill
\begin{subfigure}{0.48\textwidth}
    \includegraphics[width=\linewidth]{pictures/Anomaly_Type_Cluster.png} 
    \caption{Anomaly Types: Cluster Type AUC-ROC performance.}
    \label{fig:type_cluster}
\end{subfigure}

\vspace{0.4em}
\begin{subfigure}{0.48\textwidth}
    \includegraphics[width=\linewidth]{pictures/Anomaly_Type_Global.png} 
    \caption{Anomaly Types: Global Type AUC-ROC performance.}
    \label{fig:type_global}
\end{subfigure}\hfill
\begin{subfigure}{0.48\textwidth}
    \includegraphics[width=\linewidth]{pictures/Anomaly_Type_Local.png} 
    \caption{Anomaly Types: Local Type AUC-ROC performance.}
    \label{fig:type_local}
\end{subfigure}
\caption{AUC-ROC performance comparison of XXX and the top six baseline methods across different noise levels and anomaly types.}
\label{fig:noise_vs_tp}
\end{figure*}

For further details on rankings, parameter configurations, and additional experimental results, please see Section~\ref{exp:more_detail}.

\subsubsection{Computational Efficiency and Memory Trade-offs (RQ3)}
\label{ssubsec:RQ3}

The experimental results highlight the exceptional runtime performance of the XXX algorithm, which significantly outperforms the average runtime of other baseline methods in both training and testing phases. This efficiency is further supported by its scalability, as demonstrated in Table~\ref{tab:scale_train_time} and Table~\ref{tab:scale_infer_time}, which present the training and inference runtime performance (in seconds) of the algorithm across varying numbers of samples (\( N \)) and dimensions (\( d \)). The algorithm maintains efficient execution times even as the dataset size and dimensionality increase substantially. However, limitations are observed when handling extremely large datasets, such as \( d = 10,000 \) and \( N = 1,000,000 \), where an \textbf{Out of Memory (OoM)} error occurs. This is attributed to the memory constraints of the evaluation machine, equipped with only 64 GB of RAM, while the dataset demands over 40 GB of memory.

\begin{table}[!ht]
\caption{Training runtime performance (in seconds) of LOC-NFST for different values of \( N \) and \( d \).}
\centering
\resizebox{\linewidth}{!}{%
\begin{tabular}{|c|c|c|c|c|}
\hline
\( N \backslash d \) & \( 10^{1} \) & \( 10^{2} \) & \( 10^{3} \) & \( 10^{4} \) \\
\hline
\( 10^{2} \) & \( 0.2104 \) & \( 0.6167 \) & \( 2.1517 \) & \( 7.2927 \) \\
\( 10^{3} \) & \( 0.6018 \) & \( 2.1875 \) & \( 20.3309 \) & \( 59.2073 \) \\
\( 10^{4} \) & \( 3.2077 \) & \( 4.3860 \) & \( 63.5863 \) & \( 1.09 \times 10^{3} \) \\
\( 10^{5} \) & \( 4.2051 \) & \( 13.4534 \) & \( 336.0482 \) & \( 5.49 \times 10^{3} \) \\
\( 10^{6} \) & \( 10.5647 \) & \( 99.4214 \) & \( 3.20 \times 10^{3} \) & OoM \\
\hline
\end{tabular}%
}
\label{tab:scale_train_time}
\end{table}

\begin{table}[!ht]
\caption{Total inference runtime performance (in seconds) of LOC-NFST for different values of \( N \) and \( d \).}
\centering
\resizebox{\linewidth}{!}{%
\begin{tabular}{|c|c|c|c|c|}
\hline
\( N \backslash d \) & \( 10^{1} \) & \( 10^{2} \) & \( 10^{3} \) & \( 10^{4} \) \\
\hline
\( 10^{2} \) & \( 0.0107 \) & \( 2.49 \times 10^{-4} \) & \( 3.85 \times 10^{-4} \) & \( 0.0120 \) \\
\( 10^{3} \) & \( 0.0167 \) & \( 1.75 \times 10^{-4} \) & \( 9.91 \times 10^{-3} \) & \( 0.0232 \) \\
\( 10^{4} \) & \( 0.0122 \) & \( 5.89 \times 10^{-4} \) & \( 0.0109 \) & \( 0.3646 \) \\
\( 10^{5} \) & \( 0.0229 \) & \( 0.0280 \) & \( 0.1576 \) & \( 1.3809 \) \\
\( 10^{6} \) & \( 0.0215 \) & \( 0.0288 \) & \( 0.2306 \) & OoM \\
\hline
\end{tabular}%
}
\label{tab:scale_infer_time}
\end{table}

Figure~\ref{fig:radar_chart} provides a visual comparison of the performance metrics between the proposed model (XXX) and the baseline methods, focusing on standardized training and testing ranks, as well as train and test times, with the radar chart illustrating key insights: the blue line representing XXX shows a rank of 1.50 for Train Rank and 1.50 for Test Rank, significantly lower than the median baseline performance (orange shaded area with ranks around 6-7), indicating superior performance in both training and testing phases; additionally, the Train Time and Test Time for XXX fall within the inner shaded region, suggesting faster training and testing times compared to the median baseline, which extends further outward, collectively highlighting XXX's efficiency and accuracy advantages. The radar chart uses a normalized scale from 0 to 1, where values closer to the center indicate better performance (lower rank or shorter time), and the blue line (XXX) consistently remains inside the orange shaded area (Median of Baselines), underscoring XXX's overall advantage in terms of both efficiency and effectiveness across all evaluated metrics, reinforcing the algorithm's capability to handle large-scale data efficiently while maintaining high accuracy.

\begin{figure}[!ht]
\centering  
\includegraphics[width=\columnwidth]{pictures/radar_chart.png} 
\caption{Comparison of standardized training and testing ranks, and train/test times between XXX and baseline methods. The blue line represents XXX, while the orange shaded area indicates the median performance of baselines.}
\label{fig:radar_chart}
\end{figure}

Furthermore, we explicitly evaluate the memory and computational trade-offs of the XXX algorithm. Figure~\ref{fig:memory_footprint} visualizes the peak RAM usage during training for both XXX and the baseline methods.

\begin{figure}[!ht]
\centering
\includegraphics[width=\columnwidth]{pictures/Bar_Memory_Footprint.png} 
\caption{Comparison of Peak RAM usage (MB) during training across various models, highlighting the memory efficiency of XXX.}
\label{fig:memory_footprint}
\end{figure}

To further illustrate the balance between detection performance and memory overhead, Figure~\ref{fig:efficiency_tradeoff} presents an efficiency trade-off scatter plot. The plot contrasts the AUC-PR scores against the peak RAM consumption, demonstrating that XXX achieves state-of-the-art detection accuracy while maintaining a constrained memory footprint, making it highly suitable for IoT edge deployments.

\begin{figure}[!ht]
\centering
\includegraphics[width=\columnwidth]{pictures/Scatter_Efficiency_Tradeoff.png} 
\caption{Efficiency vs. Performance Trade-off: AUC-PR vs. Peak RAM usage.}
\label{fig:efficiency_tradeoff}
\end{figure}

Finally, Figure~\ref{fig:time_complexity} compares the computational time complexity (Train and Test Time) on a logarithmic scale, reassuring that XXX remains competitive in runtime speed against standard lightweight baselines while significantly outperforming deep learning architectures.

\begin{figure}[!ht]
\centering
\includegraphics[width=\columnwidth]{pictures/Bar_Time_Complexity.png} 
\caption{Train and Test Time complexity comparison on a logarithmic scale.}
\label{fig:time_complexity}
\end{figure}




\newpage


\subsection{Ablation Study}

We study two design choices in XXX: the number of pseudo-classes used to construct surrogate class structure, and the basis-construction method used in the projection stage. Detailed curves and per-dataset values are provided in Appendix~\emph{[APPENDIX LABEL]}.

\subsubsection{Sensitivity to the Number of Pseudo-Classes}
\label{subsec:cluster_sensitivity}

Because XXX requires at least two pseudo-classes, we vary the number of K-means clusters only over the valid range \(c \in [2,300]\). The default value used in the main experiments is \(c=15\)\, selected according to the protocol in Section~\ref{sec:exp_su}. We evaluate each setting on the four IoT datasets using AUC-PR and AUC-ROC.

\begin{figure}[!h]
\centering
\begin{subfigure}{0.49\columnwidth}
\includegraphics[width=\linewidth]{pictures/aucpr_vs_clusters.png}
\caption{AUC-PR.}
\label{fig:aucpr_vs_clusters}
\end{subfigure}%
\hfill
\begin{subfigure}{0.49\columnwidth}
\includegraphics[width=\linewidth]{pictures/aucroc_vs_clusters.png}
\caption{AUC-ROC.}
\label{fig:aucroc_vs_clusters}
\end{subfigure}
\caption{Sensitivity of XXX to the number of pseudo-classes across the four IoT datasets.}
\label{fig:au_vs_cluster}
\end{figure}

Figure~\ref{fig:au_vs_cluster} shows that performance remains relatively stable across a broad range of valid pseudo-class granularities. In particular, both AUC-PR and AUC-ROC vary only moderately as \(c\) increases, indicating that XXX is not overly sensitive to the exact number of pseudo-classes within the tested range. This is desirable in practice because the number of latent benign modes is rarely known in advance. At the same time, \(c\) still affects the geometry of the projected space and the computational cost, so it remains a tunable design parameter rather than an irrelevant one.

\subsubsection{Gram--Schmidt vs.\ SVD Basis Construction}
\label{subsec:svd_vs_gs_ablation}

We next isolate the effect of the basis-construction method used in the projection stage by comparing the SVD-based implementation with a classical Gram--Schmidt implementation. To keep this comparison focused on the projection step itself, we fix the pseudo-class count to the minimal valid setting \(c=2\). This choice is used here only to isolate the orthogonalization strategy; it is not intended to imply that \(c=2\) is the operating point used in all main experiments.

\paragraph{Detection performance}
Figure~\ref{fig:svd_gs_performance} shows that the SVD- and Gram--Schmidt-based variants achieve nearly identical AUC-PR and AUC-ROC values across datasets. This indicates that the SVD reformulation preserves the discriminative behavior of the projection stage rather than changing the scoring rule itself.

\begin{figure}[!ht]
\centering
\includegraphics[width=\columnwidth]{pictures/SVD_vs_GS_Performance.png}
\caption{AUC-PR and AUC-ROC comparison between the SVD- and Gram--Schmidt-based variants of XXX.}
\label{fig:svd_gs_performance}
\end{figure}

\paragraph{Computational efficiency}
Figure~\ref{fig:svd_gs_time} compares the computational cost of the two variants. The SVD-based implementation is consistently faster in the tested settings, which supports its use as the default basis-construction method in XXX.

\begin{figure}[!ht]
\centering
\includegraphics[width=\columnwidth]{pictures/SVD_vs_GS_Time.png}
\caption{Training-time comparison between the SVD- and Gram--Schmidt-based variants of XXX.}
\label{fig:svd_gs_time}
\end{figure}

Overall, this ablation shows that replacing Gram--Schmidt with SVD preserves detection performance while improving the practicality of projection learning in large tabular settings.











% ========================
\section{Conclusion}

This paper presented \textbf{XXX}, a lightweight benign-only novelty-detection framework for IoT intrusion detection based on a reformulation of the Null Foley--Sammon Transform (NFST). The work was motivated by three practical challenges in open-world IoT security: slight contamination in the benign training set, multimodal benign traffic together with heterogeneous anomaly manifestations, and the need for efficient deployment on large tabular network data.

To address these challenges, XXX introduces pseudo-class construction to restore surrogate discriminative structure in benign-only training, combines it with an SVD-based NFST reformulation for stable projection learning, and performs representative-based novelty scoring with a global quantile threshold in the projected space. This design preserves the lightweight nature of linear projection methods while adapting NFST to the benign-only IoT intrusion detection setting.
Experimental results on four public IoT benchmarks, together with controlled evaluations under local, global, and cluster anomaly structures and benign-set contamination levels of \(1\%\), \(3\%\), and \(5\%\), show that XXX achieves strong and consistent performance across a range of settings. Additional scalability experiments further indicate that the proposed framework remains practical for large tabular datasets in terms of runtime and memory usage.

At the same time, this study has several limitations. The current scoring rule is representative-based and does not explicitly model within-cluster covariance or local density, and the performance of the method still depends on design choices such as the number of pseudo-classes and the threshold parameter. Future work will therefore investigate more adaptive pseudo-class construction, covariance-aware or density-aware scoring in the projected space, and more systematic threshold-selection strategies for contaminated benign training data. Another promising direction is to extend the framework to streaming or continually evolving IoT environments, where benign behavior and attack patterns may drift over time.

% ========================


\bibliographystyle{elsarticle-num}
\bibliography{references}


\newpage
\appendix
\onecolumn


\section{Detailed Supplementary Theory}
\label{ap1:supp_theory}

This section provides additional theoretical insights to clarify the technical and methodological aspects related to the presented experiments. It focuses on explaining core concepts, model approaches, and characteristics of the IoT datasets used.

\subsection{Datasets}
\label{ap1:datasets}

 In this study, we thoroughly evaluate our proposed method
 on four widely used datasets in the attack detection field: BoT
IoT, CIC-IoT2023, N-BaIoT, and ToN-IoT. These datasets are
 specifically designed to simulate IoT network environments and
 include a range of up-to-date attacks. A summary of each
 dataset is provided in Table 3, and further details are described
 below:
\begin{table}[h]
\centering
\caption{Summary of IoT Datasets Used for Evaluation}
\label{tab:datasets}
\begin{tabular}{|l|c|c|c|}
\hline
\textbf{Dataset}   & \textbf{Classes} & \textbf{Features} \\
\hline
CICIoT2023         & 6                & 46                       \\
N-BaIoT            & 9                & 115                  \\
BoTIoT             & 6                & 21                \\
ToNIoT             & 10               & 21                    \\
\hline
\end{tabular}
\end{table}

 
\textbf{CICIoT2023} : This Dataset developed by the Canadian Institute for Cybersecurity, CICIoT2023 is a comprehensive dataset designed for IoT intrusion detection research~\cite{neto2023ciciot2023}. It includes traffic from various smart devices such as thermostats and cameras under both normal operations and several attack scenarios (e.g., DoS, DDoS, botnets, and reconnaissance). The dataset offers 46 network flow features, including packet sizes, protocol types, and durations. With a large sample size and diverse attack types, CICIoT2023 enables robust assessment of OC-NFST in detecting various high-dimensional and noisy anomalies.



\textbf{N-BaIoT} : The dataset targets anomaly detection in smart home environments~\cite{8490192}. It contains traffic from connected devices such as cameras and sensors under both benign and malicious conditions. Comprising 115 features—including packet counts, byte sizes, and connection timings—the dataset is particularly useful for testing clustering-based and noise-tolerant methods. Noise injection simulates real-world imperfections, testing OC-NFST’s effectiveness in identifying both global and local anomalies.


\textbf{BoTIoT}: This datasets is specifically curated for evaluating botnet detection in IoT ecosystems~\cite{koroniotis2019towards}. It includes time-series data from smart appliances and connected vehicles under normal usage and botnet-induced attacks such as scanning and DDoS. Featuring 21 attributes like packet lengths and inter-arrival times, the dataset focuses on time-dependent anomalies, making it ideal for evaluating OC-NFST’s performance on dependency-based detection in dynamic IoT scenarios.



\textbf{ToNIoT}~\cite{fesz-dm97-19}: This dataset was provided by the Cyber Range Lab of UNSW Canberra to support the development of artificial intelligence-based security solutions for industrial IoT networks. It consists of benign network traffic generated by legitimate IoT devices under normal operating conditions. Moreover, the dataset includes diverse IoT applications and communication protocols commonly used in IoT deployments.

\vspace{0.5em}


\subsection{Additional Demonstration of Synthetic Anomalies}
\label{ap1:ano_types}

To evaluate the performance of our approach, we test it on 47 core benchmark datasets, which is similar to the datasets used in the original ADBench paper \cite{han2022adbench}. This allows us to assess how well the proposed methods work across a wide range of real-world scenarios and anomaly types. Below, we describe the generation process for four common types of anomalies.

\textbf{Definition and Generation Process of Four Common Types of Anomalies} \cite{han2022adbench}:

\begin{itemize}
    \item \textbf{Local Anomalies:} Local anomalies are points that deviate significantly from their local neighborhoods. To generate local anomalies, we first create synthetic normal samples using the GMM \cite{GMM}. Then, we scale the covariance matrix by a scaling parameter $\alpha = 5$ to generate local anomalies, which alters the covariance structure of the normal data.
    
    \item \textbf{Global Anomalies:}  Global anomalies are data points that are significantly different from the normal data. These anomalies are generated using a uniform distribution $\text{Unif}(min(X_k), max(X_k))$, where the boundaries are defined by the minimum and maximum values of an input feature $X_k$. The scaling parameter $\alpha = 1.1$ controls how extreme these anomalies are.
    
   

    \item \textbf{Cluster Anomalies:} Cluster anomalies, also known as group anomalies, exhibit similar characteristics to clusters. To generate cluster anomalies, we adjust the mean feature vector of the normal samples by a scaling parameter $\alpha = 5$, i.e., shifting the mean of the GMM \cite{GMM}. This parameter controls the distance between anomaly clusters and the normal data. Then, the anomalies are generated by applying the scaled GMM \cite{GMM}.
\end{itemize}



Building on the approach from ADBench, we generate "realistic" synthetic data using a generative model like Gaussian Mixture Model (GMM) \cite{GMM}, trained on the normal samples of a benchmark dataset. The original anomalies are discarded since their exact types are unknown. By adjusting the parameters of the generative model, we can create both normal samples and different types of anomalies. This method adds complexity to the generated anomalies, making the synthetic data better reflect the challenges of real-world anomaly detection tasks.

\subsection{Baseline Algorithms}
\label{ap1:ap_baselines}

In this section, we present the baseline algorithms employed for anomaly detection. These algorithms are categorized into two groups: \textit{Shallow Machine Learning-based} and \textit{Deep Learning-based} outlier detection algorithms. 

\subsubsection{Machine Learning-based Outlier Detection Algorithms}

\begin{itemize}
    \item \textbf{Clustering-Based Local Outlier Factor (CBLOF)} \cite{CBLOF}: CBLOF identifies anomalies by first partitioning the data into clusters and then evaluating the distance between the data points and the cluster centroids. Instances that exhibit significant distance from their respective clusters are considered outliers.

    \item \textbf{K-Nearest Neighbors (KNN)} \cite{KNN}: KNN calculates anomaly scores based on the distance to the $k$-th nearest neighbor. Data points with larger distances from their nearest neighbors are assigned higher anomaly scores, suggesting that they are potential outliers.

    \item \textbf{Isolation Forest (IForest)} \cite{IForest}: Isolation Forest isolates observations by recursively partitioning the data using random features and split values. The anomalies, which are few and distinct, are isolated more efficiently, resulting in shorter path lengths.

    \item \textbf{One-Class Support Vector Machine (OCSVM)} \cite{OCSVM}: OCSVM constructs a decision boundary in a high-dimensional space, separating normal data points from outliers by maximizing the margin between the origin and the normal data points. The decision boundary thus defines the region where anomalies are unlikely to reside.

    \item \textbf{Local Outlier Factor (LOF)} \cite{LOF}: LOF measures the local density deviation of a data point in relation to its neighbors. Points exhibiting substantially lower density compared to their neighbors are considered anomalies.

    \item \textbf{Histogram-Based Outlier Detection (HBOS)} \cite{HBOS}: HBOS assumes the independence of the features and constructs histograms for each feature. Anomalies are identified based on the deviation of a data point from the expected distribution represented by these histograms.

    \item \textbf{Principal Component Analysis (PCA)} \cite{PCA}: PCA is employed to reduce the dimensionality of the data by projecting it onto a lower-dimensional space. Anomalies are detected by evaluating the reconstruction error, with larger errors indicating a poor fit to the reduced space and suggesting outlier behavior.

    \item \textbf{Subspace Outlier Detection (SOD)} \cite{SOD}: SOD detects outliers by examining various subspaces of the feature space. Anomalies are identified based on their deviation from expected behavior in one or more of these subspaces.

    \item \textbf{Connectivity-Based Outlier Factor (COF)} \cite{COF}: COF measures the chaining distance of a data point and its k-th nearest neighbors. Data points with unusually large or small chaining distances relative to their neighbors are classified as anomalies.

    \item \textbf{Lightweight Online Detector of Anomalies (LODA)} \cite{LODA}: LODA is an ensemble-based anomaly detection algorithm particularly suitable for real-time anomaly detection in streaming data. It is capable of adapting to concept drift and large-scale data environments.

    \item \textbf{Empirical Cumulative Distribution-Based Outlier Detection (ECOD)} \cite{ECOD}: ECOD uses the empirical cumulative distribution function (ECDF) to model the distribution of each feature. Anomalies are detected by identifying data points that lie in the extreme tails of the distribution.

    \item \textbf{Copula-Based Outlier Detector (COPOD)} \cite{COPOD}: COPOD is a model-free, highly interpretable anomaly detection technique that uses copula-based models to capture the dependence structure between features. Detects anomalies in multivariate data without requiring hyperparameter tuning.
\end{itemize}

\subsubsection{Deep Learning-based Outlier Detection Algorithms}

\begin{itemize}
    \item \textbf{Deep Support Vector Data Description (DeepSVDD)} \cite{DeepSVDD}: DeepSVDD uses a deep neural network to learn representations of the data while minimizing the volume of a hypersphere that encloses these representations. Data points that lie far from the center of this hypersphere are identified as anomalies.

    \item \textbf{AutoEncoder} \cite{Autoencoder}: AutoEncoders are neural networks designed to learn efficient representations of data through a process of encoding and decoding. Anomalies are detected based on reconstruction errors, with higher errors indicating instances that do not conform to the learned data distribution.

    \item \textbf{Deviation Networks (DevNet)} \cite{DevNet}: DevNet is a deep learning model that enforces a deviation-based loss function to penalize large deviations in the learned representations. This method is particularly effective in identifying outliers by ensuring that anomalous instances exhibit substantial deviations from the learned data distribution.

    \item \textbf{Learning Unsupervised Anomalies with Neural Representation learning (LUNAR)} \cite{LUNAR}: LUNAR is a deep learning-based, neural network approach designed for unsupervised anomaly detection. It learns a data distribution through representation learning and detects anomalies by identifying instances that deviate significantly from the normal patterns learned. This method is particularly effective in high-dimensional, complex data environments where labeled data is unavailable.



    \item \textbf{Autoencoder-based One-Class Support Vector Machine (AE1SVM)} \cite{AE1SVM}: AE1SVM is a hybrid method combining AutoEncoder-based feature learning with Support Vector Machines for enhanced anomaly detection. This approach leverages the reconstruction errors from the AutoEncoder in conjunction with SVM’s classification capabilities to detect outliers.

    \item \textbf{Adversarial Learned Anomaly Detection (ALAD)} \cite{ALAD}: ALAD (additional details to be included) is a deep learning model that utilizes adversarial learning techniques for anomaly detection. The architecture is designed to capture the subtle differences between normal and anomalous instances by learning the underlying data distribution.
\end{itemize}

\subsection{Evaluation Metrics}
\label{ap1:EM}
To comprehensively assess the performance of the proposed method, we adopt a suite of evaluation metrics that capture various aspects of binary classification performance, especially relevant to anomaly detection in imbalanced and noisy IoT network traffic. The formulas and interpretations are as follows:

\begin{itemize}
    \item \textbf{Accuracy (ACC):} Measures the proportion of correctly predicted instances among all samples.
    \[
    \text{ACC} = \frac{\text{TP} + \text{TN}}{\text{TP} + \text{TN} + \text{FP} + \text{FN}}
    \]

    \item \textbf{True Positive Rate (TPR) or Recall:} Indicates the ability to correctly identify actual anomalies.
    \[
    \text{TPR} = \frac{\text{TP}}{\text{TP} + \text{FN}}
    \]

    \item \textbf{False Positive Rate (FPR):} Represents the proportion of normal instances incorrectly classified as anomalies.
    \[
    \text{FPR} = \frac{\text{FP}}{\text{FP} + \text{TN}}
    \]

    \item \textbf{Positive Predictive Value (PPV) or Precision:} Measures how many of the instances predicted as anomalies are actually anomalous.
    \[
    \text{PPV} = \frac{\text{TP}}{\text{TP} + \text{FP}}
    \]

    \item \textbf{F1-score:} Harmonic mean of precision and recall, providing a balanced metric for imbalanced datasets.
    \[
    \text{F1-score} = \frac{2 \cdot \text{TPR} \cdot \text{PPV}}{\text{TPR} + \text{PPV}}
    \]

    \item \textbf{Matthews Correlation Coefficient (MCC):} A correlation-based metric that considers all four values in the confusion matrix and is robust to imbalanced datasets.
    \[
    \text{MCC} = \frac{\text{TP} \cdot \text{TN} - \text{FP} \cdot \text{FN}}{\sqrt{(\text{TP} + \text{FP})(\text{TP} + \text{FN})(\text{TN} + \text{FP})(\text{TN} + \text{FN})}}
    \]

    \item \textbf{Area Under the ROC Curve (AUC-ROC):} Quantifies the model’s ability to distinguish between classes across various thresholds. It is computed as the area under the ROC curve, which plots TPR versus FPR.
    \[
    \text{AUC-ROC} = \sum_{i=1}^{n-1} \frac{\text{TPR}_{i+1} + \text{TPR}_i}{2} \cdot (\text{FPR}_{i+1} - \text{FPR}_i)
    \]
    A value closer to 1 indicates stronger class separation capability.

    \item \textbf{Area Under the Precision-Recall Curve (AUC-PR):} Captures the trade-off between precision and recall, especially informative in scenarios with significant class imbalance. It is typically approximated using the trapezoidal rule:
    \[
    \text{AUC-PR} = \sum_{i=1}^{n-1} (\text{Recall}_{i+1} - \text{Recall}_i) \cdot \frac{\text{Precision}_{i+1} + \text{Precision}_i}{2}
    \]
    Higher AUC-PR values reflect better performance in identifying positive (anomalous) instances.
\end{itemize}

This comprehensive set of metrics provides a robust foundation for evaluating the accuracy, reliability, and robustness of our proposed method under diverse and challenging IoT traffic scenarios.


\section{Experiments with More Details and Results}
\label{ap2:experiments}

This section details the experimental setup and results, analyzing model performance across IoT datasets under varying noise levels and anomaly types. Accuracy and AUCROC scores, reported in Tables \ref{tab:model_noise_perf} and \ref{tab:aucroc_results}, highlight OurModel’s superiority, achieving up to 96.29\% accuracy on BIoT (1\% noise) and 100\% AUCROC for global anomalies. A survey of recent IDS research in Table \ref{tab:ids_overview} compares 11 studies, showing diverse approaches with accuracies from 45\% to 99\%, contextualizing our method’s advancements.


\subsection{Expriments results with more details}
\label{ap2:more_detail}
This section presents a detailed evaluation of various models across multiple IoT datasets under different noise levels and anomaly types. Table \ref{tab:model_noise_perf} reports the model performance under varying noise conditions (1\%, 3\%, 5\%) on BIoT, CICIoT2023, NBaIoT, and ToNIoT datasets, highlighting the robustness of each model. Table \ref{tab:aucroc_results} provides AUCROC scores across local, global, and cluster anomaly types, showcasing the models' effectiveness in detecting different anomaly patterns.


Table \ref{tab:model_noise_perf} summarizes the performance of various models under different noise levels (1\%, 3\%, 5\%) across four IoT datasets: BIoT, CICIoT2023, NBaIoT, and ToNIoT. OurModel consistently outperforms others, achieving the highest scores across all noise levels and datasets, with top accuracies like 96.29\% (BIoT, 1\% noise) and 89.53\% (ToNIoT, 1\% noise). KNN follows closely, particularly strong in BIoT and ToNIoT, while AutoEncoder and CBLOF show competitive results in specific scenarios. Other models, such as DevNet and COPOD, exhibit lower performance, especially under higher noise levels. 

\begin{table*}[!ht]
\caption{Model performance across different noise levels (1\%, 3\%, 5\%) on various IoT datasets.}
\centering
\resizebox{\textwidth}{!}{%
\begin{tabular}{|c|ccc|ccc|ccc|ccc|}
\hline
\textbf{Model} & \multicolumn{3}{c|}{\textbf{BIoT}} & \multicolumn{3}{c|}{\textbf{CICIoT2023}} & \multicolumn{3}{c|}{\textbf{NBaIoT}} & \multicolumn{3}{c|}{\textbf{ToNIoT}} \\
\cline{2-13}
 & 1\% & 3\% & 5\% & 1\% & 3\% & 5\% & 1\% & 3\% & 5\% & 1\% & 3\% & 5\% \\
\hline
OurModel & \textcolor{red}{96.29} & \textcolor{red}{93.77} & \textcolor{red}{92.90} & \textcolor{red}{95.37} & \textcolor{red}{93.93} & \textcolor{red}{93.35} & \textcolor{red}{79.24} & \textcolor{red}{77.36} & \textcolor{red}{75.35} & \textcolor{red}{89.53} & \textcolor{red}{83.92} & \textcolor{red}{81.47} \\
KNN & \textcolor{blue}{96.00} & \textcolor{blue}{90.63} & \textcolor{blue}{88.24} & \textcolor{blue}{93.11} & \textcolor{blue}{92.43} & 90.09 & 78.15 & \textcolor{blue}{76.02} & \textcolor{blue}{74.81} & \textcolor{blue}{86.31} & \textcolor{blue}{83.54} & \textcolor{blue}{79.79} \\
AutoEncoder & 95.66 & 89.62 & 86.45 & 92.46 & 91.64 & \textcolor{blue}{ 91.24}& 55.89 & 21.90 & 66.72 & 66.32 & 67.32 & 71.73 \\
CBLOF & 87.06 & 86.50 & 85.99 & 90.68 & 90.65 & 89.90 & 75.28 & 75.88 & 74.70 & 74.35 & 73.92 & 76.69 \\
AE1SVM & 92.43 & 91.41 & 86.94 & 90.19 & 90.77 & 90.44 & 56.06 & 64.31 & 48.93 & 54.71 & 56.57 & 52.80 \\
DevNet & 64.14 & 63.63 & 64.11 & 85.86 & 85.82 & 85.85 & \textcolor{blue}{ 78.28} & 64.49 & 64.49 & 76.31 & 76.36 & 76.30 \\
PCA & 90.19 & 88.21 & 86.07 & 89.97 & 89.44 & 88.62 & 69.16 & 66.30 & 61.66 & 53.64 & 52.46 & 52.10 \\
IForest & 84.24 & 80.56 & 82.09 & 92.29 & 92.18 & 91.66 & 30.87 & 29.73 & 27.18 & 76.57 & 75.38 & 75.80 \\
LODA & 76.00 & 91.55 & 79.47 & 83.50 & 81.36 & 86.99 & 47.32 & 41.75 & 63.17 & 56.56 & 69.65 & 62.53 \\
DeepSVDD & 87.27 & 93.63 & 79.28 & 86.40 & 90.33 & 89.47 & 62.17 & 61.35 & 57.68 & 63.10 & 62.24 & 63.24 \\
HBOS & 85.71 & 84.57 & 83.79 & 94.13 & 80.96 & 78.59 & 31.01 & 29.72 & 28.84 & 82.05 & 81.49 & 80.74 \\
LOF & 92.38 & 65.33 & 61.21 & 88.92 & 82.26 & 75.06 & 76.46 & 57.45 & 56.32 & 82.68 & 76.36 & 70.72 \\
MO\_GAAL & 39.89 & 44.38 & 45.17 & 54.96 & 55.10 & 55.68 & 59.74 & 31.31 & 35.14 & 62.32 & 77.88 & 77.78 \\
OCSVM & 85.75 & 83.69 & 81.42 & 84.39 & 84.01 & 83.40 & 77.63 & 76.66 & 30.47 & 48.38 & 48.11 & 48.31 \\
ALAD & 53.62 & 53.10 & 49.96 & 50.57 & 51.17 & 51.14 & 43.33 & 43.30 & 45.49 & 51.90 & 51.52 & 52.41 \\
COPOD & 40.73 & 40.66 & 40.63 & 44.90 & 44.88 & 44.86 & 31.88 & 31.77 & 31.61 & 29.81 & 29.79 & 29.75 \\
ECOD & 42.34 & 42.27 & 42.24 & 34.86 & 34.86 & 34.84 & 30.44 & 30.25 & 30.11 & 23.59 & 23.57 & 23.55 \\
\hline
\end{tabular}
}
\label{tab:model_noise_perf}
\end{table*}

Table \ref{tab:aucroc_results} presents AUCROC scores for models across local, global, and cluster anomaly types on the same IoT datasets. OurModel again leads, achieving near-perfect scores (e.g., 100\% for global and cluster anomalies in BIoT, NBaIoT, and ToNIoT). KNN, AutoEncoder, and PCA show strong performance, particularly for global and cluster anomalies, with scores often exceeding 99\%. Models like DevNet and ALAD struggle with local anomalies, while IForest and HBOS perform well in specific cases, such as global anomalies in BIoT and NBaIoT.


\begin{table}[!ht]
\caption{AUCROC Scores for Different Models across Anomaly Types and Datasets}
\centering
\resizebox{\textwidth}{!}{%
\begin{tabular}{|c|ccc|ccc|ccc|ccc|}
\hline
\textbf{Model} & \multicolumn{3}{c|}{\textbf{BoT-IoT}} & \multicolumn{3}{c|}{\textbf{CICIoT2023}} & \multicolumn{3}{c|}{\textbf{N-BaIoT}} & \multicolumn{3}{c|}{\textbf{ToN-IoT}} \\
 & Local & Global & Cluster & Local & Global & Cluster & Local & Global & Cluster & Local & Global & Cluster \\
\hline
OurModel & \textcolor{red}{95.32} & \textcolor{red}{100} & \textcolor{red}{100} & 91.29 & \textcolor{red}{99.94} & \textcolor{red}{100} & \textcolor{red}{92.76} & \textcolor{red}{100} & \textcolor{red}{99.99} & \textcolor{red}{93.16} &\textcolor{red}{100} & \textcolor{red}{ 100} \\
KNN & 90.11 & 99.93 & 99.93 & 70.90 & 99.76 & 99.73 & 83.62 & 99.59 & \textcolor{blue}{99.85} & 85.89 & 99.93 & 99.95 \\
AutoEncoder & 89.58 & 99.83 & 99.86 & 95.44 & 99.35 & 99.56 & 89.40 & 99.92 & \textcolor{blue}{99.85} & \textcolor{blue}{91.11} & 99.89 & 99.94 \\
CBLOF & 89.49 & 99.72 & 99.93 & 69.51 & 99.69 & 99.63 & 85.07 &  \textcolor{blue}{ 99.99} & 99.70 & 85.38 & 99.93 & \textcolor{blue}{99.96} \\
AE1SVM & 90.41 & 99.80 & 99.93 & 95.16 & 99.56 &  \textcolor{blue}{99.78} & 89.53 & 99.85 & 99.77 & 89.03 & 99.93 & 99.88 \\
DevNet & 52.71 & 24.49 & 3.48 & 54.50 & 54.29 & 11.40 & 51.38 & 41.33 & 3.40 & 52.05 & 31.88 & 31.19 \\
PCA & 87.26 & 99.66 & 99.86 & \textcolor{blue}{95.24} &  \textcolor{blue}{99.78} & \textcolor{blue}{99.78} & 86.34 & \textcolor{red}{100} & 99.92 & 84.93 & 99.88 & 99.88 \\
IForest & 83.23 & \textcolor{blue}{99.99} & \textcolor{red}{100} & 94.34 & 99.33 & 99.56 & 92.72 & \textcolor{red}{100} & 99.78 & 84.38 & \textcolor{red}{100} & 99.75 \\
LODA & 89.57 & 99.45 & 99.84 & 82.10 & 96.80 & 99.52 & \textcolor{blue}{92.56} & 99.27 & 99.93 & 76.45 & 99.96 & 98.66 \\
DeepSVDD & 86.47 & 99.58 & 99.93 & 90.49 & 99.55 & 99.55 & 90.20 & 99.91 & 99.93 & 91.92 & 99.92 & 99.94 \\
HBOS & 77.33 & 99.93 &   \textcolor{red}{100} & 90.37 & 99.55 & 96.98 & 92.70 & \textcolor{red}{100} & 99.99 & 79.38 & 99.94 & 99.88 \\
LOF & \textcolor{blue}{92.63} & 99.89 & 99.93 & 90.71 & 99.76 & 99.86 & 88.65 & 99.54 & 99.70 & 90.59 & 99.87 & \textcolor{red}{100} \\
MO\_GAAL & 55.83 & 99.85 & 99.70 & 65.05 & 66.52 & 99.68 & 62.35 & 67.76 & 99.68 & 65.07 & 57.43 & 97.56 \\
OCSVM & 90.56 & 99.72 & 99.93 & 89.85 & 99.28 & 99.77 & 92.33 & 99.45 & 99.63 & 88.32 & 99.94 & 99.98 \\
ALAD & 54.34 & 61.86 & 83.01 & 51.74 & 60.39 & 86.63 & 50.65 & 57.89 & 89.39 & 55.67 & 62.02 & 75.87 \\
COPOD & 73.89 & 99.76 & 30.07 & 92.45 & 99.73 & 56.28 & 91.25 & \textcolor{red}{100} & 12.68 & 92.76 & \textcolor{blue}{99.99} & 54.17 \\
ECOD & 78.26 & 99.82 & 48.26 & \textcolor{red}{96.59} & 99.91 & 44.45 & 85.69 & \textcolor{blue}{99.99} & 19.63 & 92.22 & \textcolor{blue}{99.99} & 50.47 \\
\hline
\end{tabular}%
}
\label{tab:aucroc_results}
\end{table}


\subsection{Survey of Competitor Results}
\label{ap2:competitor_}

This section provides an overview of recent research on Intrusion Detection Systems (IDSs) for IoT networks, as summarized in Table \ref{tab:ids_overview}. The table compiles key details from 11 studies, including their titles, years, methodologies, datasets, tasks, models, and performance metrics. The surveyed works primarily focus on anomaly detection, binary classification, and novelty detection, utilizing datasets such as BoTIoT, CICIoT2023, N-BaIoT, and TON-IoT. Notable models include DeepSVDD, OCSVM, KANs, and RNN-based approaches, with performance metrics ranging widely (e.g., accuracy from 45\% to 99\%, F1 scores from 62\% to 98\%). These studies highlight diverse approaches to enhancing IoT security, with varying degrees of effectiveness across different datasets and tasks.

% ===================== Table in full-page =====================
\onecolumn

% Trong nội dung:
\begin{table*}[p] % full-page width
\centering
\caption{Overview of recent research on IDSs, summarizing training approach, task, model, datasets, and performance.}
\label{tab:ids_overview}
\renewcommand{\arraystretch}{2.1} % tăng khoảng cách dòng
% \small  % chữ to hơn nếu cần \normalsize hoặc \large
\scriptsize % Cỡ nhỏ hơn small 
\setlength{\tabcolsep}{4pt} % giảm khoảng cách giữa cột

\begin{adjustbox}{width=1\textwidth,center} % có thể tăng lên nếu cần
\begin{tabular}{@{}c>{\centering\arraybackslash}m{3.5cm}c>{\arraybackslash}m{4.8cm}c>{\centering\arraybackslash}m{2.8cm}c>{\centering\arraybackslash}m{2.5cm}c>{\centering\arraybackslash}m{3.2cm}@{}}
\toprule
\textbf{ID} & \textbf{Title} & \textbf{Year} & \textbf{Summary} & \textbf{Dataset} & \textbf{Tasks} & \textbf{Model} & \textbf{Performance} \\ 
\midrule    

1 & Exploring Unsupervised One-Class Classifiers for Lightweight Intrusion Detection in IoT Systems\cite{golestani2024exploring}& 2024 & Compares five one-class classifiers with ML/DL models in IoT & BoTIoT & Anomaly detection & iForest + OCSVM + LOF + DeepSVDD + DROCC & ACC: 45\%-97\%, PRE: 100\%, REC: 45\%-97\%, F1: 62\%-98\% \\

2 & Device-Specific Anomaly Detection Models for IoT Systems\cite{golestani2024device} & 2024 & Compare single models and device -specific model & N-BaIoT, CICIoT2023 & Anomaly Detection & DeepSVDD + LOF + RF + DT & ACC: 84\%-99\%, F1: 88\%-99\% \\

3 & Classifying IoT Botnet Attacks With Kolmogorov-Arnold Networks: A Comparative Analysis of Architectural Variations \cite{10839389}
 & 2024 & Kolmogorov-Arnold Networks (KANs) and their variations in detecting IoT botnet attacks & N-BaIoT &  Anomaly Detection & ORIGINAL-KAN & PRE: 95\%, REC: 91\%, F1:88\%\\

4 &Robust Anomaly Detection in IoT Networks using Deep SVDD and Contractive Autoencoder\cite{aktar2024robust} & 2024 & Using Deep SVDD and Contractive Autoencoder in Anomaly Detection
 & TON-IoT &  Anomaly Detection & KMeans+OCSVM+Isolation Forest+DSVDD-CAE & ACC: 88-99\%, PRE: 84-98\%, REC: 72-99\%, F1: 76-99\% \\

% 5 Chưa xong 
5 & Firefy algorithm based WSN‑IoT security enhancement with machine learning for intrusion detection\cite{karthikeyan2024firefly} & 2024 & RNN-based IDS & TON-IoT  & Intrusion Detection OR?  Anomaly Detection  & RNN/LSTM/GRU & ACC: 94–85\% \\ 

6 &A Genetic Algorithm and t-Test-Based System for DDoS Attack Detection in IoT Networks\cite{saiyed2024genetic} & 2024 & Genetic Algorithm and t-Test-Based System for DDoS Attack Detection in IoT Networks &TON-IoT &  binary classification & GADAT-ET  & ACC: 95\%, PRE: 95\%, REC: 95\%, F1: 94\% \\

7 &Investigation of the impact effectiveness of adversarial data leakage attacks on the machine
learning models\cite{parfenov2024investigation}  & 2023 & Adversarial attacks in detecting anomalous activity in IoT networks  & CICIoT2023  & binary classification & MLP-Prod + Catboost & ACC: 95\% \\

6 &A Genetic Algorithm and t-Test-Based System for DDoS Attack Detection in IoT Networks\cite{saiyed2024genetic} & 2024 & Genetic Algorithm and t-Test-Based System for DDoS Attack Detection in IoT Networks & TON-IoT &  Binary classification & GADAT-ET  & ACC: 95\%, PRE: 95\%, REC: 95\%, F1: 94\% \\

7 &Investigation of the impact effectiveness of adversarial data leakage attacks on the machine
learning models\cite{parfenov2024investigation}  & 2023 & Adversarial attacks in detecting anomalous activity in IoT networks  & CICIoT2023  & binary classification & MLP-Prod + Catboost & ACC: 95\% \\

8 & Comparison of Multiple Feature Selection Techniques for Machine Learning-Based Detection of IoT Attacks\cite{phan2024comparison} & 2024 & Experience of feature selection methods & CICIoT2023 & Novelty Detection & K-NN + MLP & ACC:83-92\%, PRE: 82-91\%, REC: 82-93\%\\

%  9 chưa xong 
9 & Towards the Development of Realistic Botnet Dataset in the Internet of Things for Network Forensic Analytics: Bot-IoT Dataset\cite{koroniotis2019towards} & 2019 &   A new dataset & BoTIoT & Binary classification & RNN & ACC: 97\%, PRE: 100\%, REC: 97\% \\

10 & Unsupervised Anomaly Detection Approach for Cyberattack Identification\cite{segurola2024unsupervised} & 2024 & Three-step unsupervised anomaly detection (empirical Bayesian scoring, feature selection, model election) & BoTIoT & Anomaly Detection & RSR + EAD  & REC: 97\%, F1:98\%.  \\ 

11 & Improving the IDS for BoT-IoT Dataset-Based Machine Learning Classifiers \cite{9993869} & 2022 & Experiences on BoTIoT dataset & BoTIoT & Binary classification  & RF  & REC: 96\%, F1:96\%.  \\ 

\bottomrule
\end{tabular}
\end{adjustbox}
\end{table*}

\twocolumn




\subsubsection{Comparison with State-of-the-Art}
To comprehensively assess the effectiveness of our proposed method, \textbf{XXX}, we compare it against several recent state-of-the-art approaches published from 2022 onward. These methods span a range of machine learning and deep learning techniques applied to IoT intrusion detection, including optimization-based models~\cite{aktar2024robust}, genetic algorithms~\cite{saiyed2024genetic}, ensemble frameworks~\cite{golestani2024device}, and unsupervised learning techniques~\cite{segurola2024unsupervised}. The evaluation is conducted on four widely-used public datasets: \textbf{N-BaIoT}, \textbf{ToN-IoT}, \textbf{CICIoT2023}, and \textbf{BoT-IoT}.


Table~\ref{tab:ids_comparison} reports the performance of all methods in terms of standard classification metrics: \textbf{Accuracy (ACC)}, \textbf{Precision (PRE)}, \textbf{Recall (REC)}, and \textbf{F1-score (F1)}. Our method, XXX, achieves consistently competitive and often superior results across all datasets and metrics.

On the \textbf{N-BaIoT dataset}, XXX achieves a perfect score of 100\% across all four metrics. While other recent works~\cite{golestani2024device,10839389} report competitive F1-scores (e.g., 92--99\%), none surpass XXX in balanced and complete performance.

On the \textbf{ToN-IoT dataset}, XXX maintains high performance with 94\% accuracy and an F1-score of 96\%, outperforming other approaches that show broader performance variation~\cite{aktar2024robust, karthikeyan2024firefly, saiyed2024genetic}. Notably, although~\cite{saiyed2024genetic} reports similar accuracy, XXX achieves better consistency across precision and recall.

In the more challenging \textbf{CICIoT2023 dataset}, which contains complex and imbalanced traffic patterns, XXX delivers 100\% precision and a solid 92\% F1-score. While some prior works such as~\cite{parfenov2024investigation, phan2024comparison} reach comparable accuracy ranges, they lack consistency across precision and recall, and often do not report complete metric sets.

On the \textbf{BoT-IoT dataset}, XXX reaches 95\% accuracy with 99\% precision and a 97\% F1-score, closely rivaling the strongest reported results~\cite{koroniotis2019towards, golestani2024exploring}. Importantly, XXX achieves stable performance across all evaluation criteria, demonstrating robustness in the presence of noisy and imbalanced IoT data.

In summary, the results validate the \textbf{superior detection capability and generalizability} of XXX across different datasets and attack types. Its strength lies in effectively preserving local data structure, mitigating the influence of outliers, and maintaining high recall—key requirements for practical anomaly detection systems in dynamic IoT environments.
\begin{table*}[!ht]
\caption{Performance comparison of our proposed method (XXX) and recent research on four IoT datasets using Accuracy, Precision, Recall, and F1-score.}
\centering
\small % Smaller font for better scaling
\setlength{\tabcolsep}{6pt} % Increase column spacing
\resizebox{\textwidth}{!}{%
\begin{tabular}{|c|c|c|c|c|c|c|}
\hline
\textbf{Dataset} & \textbf{Method} & \textbf{ACC (\%)} & \textbf{PRE (\%)} & \textbf{REC (\%)} & \textbf{F1 (\%)} \\ 
\hline
{\textbf{N-BaIoT}} & \textbf{XXX
(Our model)} & 100.00 & 100.00 & 100.00 & 100.00 \\ 
& Ref \cite{golestani2024device} & 88.00--99.00 & \text{--} & \text{--} & 92.00--99.00 \\
& Ref \cite{10839389} & \text{--} & 95.00 & 91.00 & 88.00 \\
\hline
{\textbf{ToN-IoT}} & \textbf{XXX} & 94.00 & 98.00 & 95.00 & 96.00 \\ 
& Ref \cite{aktar2024robust} & 88.00--99.00 & 84.00--98.00 & 72.00--99.00 & 76.00--99.00 \\
& Ref \cite{karthikeyan2024firefly} & 85.00--94.00 & \text{--} & \text{--} & \text{--} \\
& Ref \cite{saiyed2024genetic} & 95.00 & 95.00 & 95.00 & 94.00 \\
\hline
{\textbf{CICIoT2023}} & \textbf{XXX} & 85.00 & 100.00 & 85.00 & 92.00 \\
& Ref \cite{parfenov2024investigation} & 95.00 & \text{--} & \text{--} & \text{--} \\
& Ref \cite{golestani2024device} & 84.00--90.00 & \text{--} & \text{--} & 88.00--92.00 \\
& Ref \cite{phan2024comparison} & 83.00--92.00 & 82.00--91.00 & 82.00--93.00 & \text{--} \\
\hline
{\textbf{BoTIoT}} & \textbf{XXX} & 95.00 & 99.00 & 95.00 & 97.00 \\
& Ref \cite{koroniotis2019towards} & 97.00 & 100.00 & 97.00 & \text{--} \\
& Ref \cite{segurola2024unsupervised} & \text{--} & \text{--} & 97.00 & 98.00 \\
& Ref \cite{9993869} & \text{--} & \text{--} & 96.00 & 96.00 \\
& Ref \cite{golestani2024exploring} & 45.00--97.00 & 100.00 & 45.00--97.00 & 62.00--98.00 \\
\hline
\end{tabular}
}
\label{tab:ids_comparison}
\end{table*}




\end{document}


% % ========================
% \section{Experiment Set Up}
% \label{sec:exp_su}
% % ========================
% To address the challenges outlined in the Introduction, this section details the experimental setup and the key research questions that guide our evaluation:

% \begin{itemize}
%     \item RQ1 (Detection Performance in IoT Networks): Does LON outperform leading anomaly detection algorithms across diverse IoT network scenarios?
    
%     \item RQ2 (Robustness Across Noise and Anomaly Types): How robustly does LON perform under complex real-world conditions, including varying noise levels and diverse anomaly types, in IoT traffic?
    
%     \item RQ3 (Runtime Efficiency and Scalability): How does LON perform in terms of computational efficiency when applied to large-scale IoT datasets?
% \end{itemize}

% % Describing datasets and their configurations
% \subsection{Datasets}
% \label{subsec:datasets}
% To address \textbf{RQ1}, we use four IoT datasets, each with 100,000 samples: \textbf{CICIoT2023} and \textbf{N-BaIoT} capture diverse network attacks and IoT device traffic for anomaly detection, while \textbf{BoTIoT} and \textbf{ToNIoT} focus on botnet activities and real-world IoT traffic, enabling comprehensive analysis of varied IoT scenarios ~\cite{neto2023ciciot2023, 8490192, koroniotis2019towards, fesz-dm97-19}. (detail described in  \ref{ap1:datasets})
% \vspace{1em}

% For answering \textbf{RQ2}, we generate four synthetic datasets derived from the four real-world IoT datasets using the definitions and methodology from \cite{han2022adbench}, incorporating local, global, and  cluster anomalies. Although cluster anomalies are included, their inter-sample dependencies are less relevant to typical IoT traffic (see more in \ref{ap1:ano_types}) . Anomalies are injected at noise levels of 1\%, 3\%, and 5\% to simulate varying degrees of contamination. This design captures the complexity of real-world IoT networks, including diverse anomaly behaviors and background noise, providing a challenging environment for evaluating detection robustness.

% To answer \textbf{RQ3}, we employ a random dataset with sample sizes ranging from 1,000 to 1,000,000 and feature dimensions ranging from 10 to 10,000.

% \vspace{1em}


% \subsection{Experiment Training and Testing Strategy}
% \label{subsec:exp_t&t}
% \textbf{Train-Test Split:} Each dataset is divided into training and testing subsets using an 80:20 ratio. Specifically, 70\% of the data is used for training and 30\% for testing across all research questions (RQ1, RQ2, and RQ3). This split ensures the model is exposed to sufficient normal traffic during training while preserving a representative and diverse test set for robust evaluation.

% \textbf{Data Cleaning:} Duplicate records are removed and missing values are handled to ensure data quality and consistency. This preprocessing step helps reduce noise and prevents the model from learning spurious patterns.

% \textbf{Feature Encoding:} Non-numeric attributes are converted into numerical representations to ensure compatibility with the model. This enables categorical features to be effectively integrated into the training process.

% \textbf{Normalization:} Feature values are normalized to ensure consistent scales across attributes, which improves training stability and facilitates faster convergence.

% \subsection{Baselines and Competitors}
% \label{subsec:bs&cp}
% We utilize the PyOD library \cite{zhao2019pyod}, a widely used benchmarking framework for anomaly detection, to implement the baseline methods. To ensure a comprehensive comparison, we select 18 state-of-the-art algorithms—including 12 traditional machine learning methods and 6 deep learning-based approaches—based on the large-scale benchmark study in \cite{han2022adbench}. These competitors represent leading anomaly detection models developed from 2020 to the present and are evaluated across all four IoT datasets. The detailed baseline algorithms are listed in \ref{ap1:ap_baselines}.  

% \subsection{Evaluation Metrics}
% \label{subsec:EM}
% To evaluate our approach, we use a set of metrics that capture detection accuracy, robustness to anomaly types and noise, and computational efficiency. These include standard classification metrics from the confusion matrix: ACC, TPR/Recall, FPR, PPV, F1-score, and MCC.

% For \textbf{RQ1} and \textbf{RQ2}, we also report AUC-ROC and AUC-PR, which are more reliable under class imbalance and varying anomaly severity. AUC-ROC evaluates class separability across thresholds, while AUC-PR emphasizes precision-recall trade-offs in imbalanced settings.

% For \textbf{RQ3}, which assesses runtime efficiency and scalability, we measure the \textbf{execution time (in seconds)} using Python’s built-in timing utilities. This allows us to capture the computational cost across datasets with varying sizes and feature dimensions.

% Detailed definitions and mathematical formulations for all metrics are provided in \ref{ap1:EM}.




% \subsection{ Software and hardware environment configuration}
% \label{subsec:ENV}
% We run our experiments in a Python 3.10.13 environment with widely adopted libraries for machine learning and data analysis. All training and evaluation tasks are performed on a server running Ubuntu 22.04, equipped with a 13th Gen Intel\textsuperscript{\textregistered} Core\textsuperscript{TM} i5-13400F processor and 32\,GB of RAM.

% Since there are no publicly available Python implementations of the original NFST algorithm, we re-implemented it based on the kernel-based extension described in its original publication~\cite{bodesheim2013kernel}. In all subsequent experiments, we refer to this implementation as \textbf{KNFST}.

% For baseline comparison, we use several state-of-the-art outlier detection methods available in the \texttt{PyOD} framework~\cite{jiang2023adgym, han2022adbench}, version 1.0.8, which provides well-documented implementations of competitive algorithms. Additionally, we use the \texttt{Scikit-learn} library~\cite{pedregosa2011scikit}, version 1.3.2, for data pre-processing and the computation of evaluation metrics.



% \section{Experiment results and competitor}
% \label{sec:exp_r&d}

% This section focuses on evaluating the efficacy of XXX, presenting a comprehensive analysis of its performance across various experimental setups to validate its effectiveness in anomaly detection within IoT environments. The evaluation leverages multiple datasets and advanced analytical techniques to provide robust insights into the method's capabilities.

% \subsection{Experiment Results}
% This section focuses on evaluating the efficacy of XXX, delving into detailed performance metrics and comparative analyses to demonstrate its superiority over baseline methods. The results are derived from rigorous experiments conducted on diverse IoT datasets, offering a solid foundation for understanding its practical applicability and potential enhancements.


% \label{subsec:exp_rs}
% \subsubsection{Anomaly Detection Performance on 4 Datasets IoT (RQ1)}
% \label{ssubsec:RQ1}
% This section comprehensively evaluates the performance of the proposed XXX method across four distinct IoT datasets—BoIoT, CICIoT2023, N-BaIoT, and ToN-IoT—using multiple metrics, as detailed in Table~\ref{tab:iot_metrics}. The evaluation relies on Area Under the Receiver Operating Characteristic Curve (AUC-ROC), Area Under the Precision-Recall Curve (AUC-PR), Matthews Correlation Coefficient (MCC), Accuracy, F1 Score, Precision, and Recall, highlighting the method's robustness. AUC-ROC scores reach 99.39\% for N-BaIoT and 97.10\% for BoIoT, indicating strong classification performance, while AUC-PR exceeds 98\%, with CICIoT2023 achieving 99.84\%, reflecting excellent precision-recall balance. MCC peaks at 0.9865 for N-BaIoT, and Accuracy hits 99.77\% for the same dataset, demonstrating high reliability. F1 scores like 0.9987 for N-BaIoT and 0.9692 for BoIoT, alongside Precision of 0.9998 for CICIoT2023 and Recall of 0.9998 for N-BaIoT, underscore a balanced and effective anomaly detection capability across diverse IoT environments.

% \begin{table*}[!ht]
% \caption{Performance of the method on four IoT datasets using AUC-ROC, AUC-PR, MCC, Accuracy, F1, Precision, and Recall.}
% \centering
% \resizebox{\textwidth}{!}{%
% \begin{tabular}{|c|c|c|c|c|c|c|c|}
% \hline
% \textbf{Dataset} & \textbf{AUC-ROC (\%)} & \textbf{AUC-PR (\%)} & \textbf{MCC} & \textbf{ACC (\%)} & \textbf{F1} & \textbf{Precision} & \textbf{Recall} \\
% \hline
% BoIoT      & 97.10 & 99.52 & 0.7378 & 94.53 & 0.9692 & 0.9891 & 0.9502 \\
% CICIoT2023 & 94.93 & 99.84 & 0.3845 & 85.82 & 0.9211 & 0.9998 & 0.8540 \\
% N-BaIoT    & 99.39 & 99.79 & 0.9865 & 99.77 & 0.9987 & 0.9977 & 0.9998 \\
% ToN-IoT    & 94.63 & 98.74 & 0.7349 & 93.76 & 0.9641 & 0.9809 & 0.9478 \\
% \hline
% \end{tabular}
% }
% \label{tab:iot_metrics}
% \end{table*}


% % Required packages in the preamble:
% % \usepackage[table]{xcolor}
% % \usepackage{multirow}
% % \usepackage{adjustbox}


% The analysis extends to clustering techniques, as visualized in Figure~\ref{fig:au_vs_cluster}, where AUC-PR and AUC-ROC are assessed using the K-means algorithm with varying cluster numbers. These metrics reveal the method's adaptability and robustness, with AUC-PR scores around 98-99\% and AUC-ROC scores approximately 95-99\%, indicating strong precision and classification accuracy. The AUC-PR plot (Figure~\ref{fig:aucpr_vs_clusters}) shows XXX maintaining near 99\% performance despite increasing clusters, crucial for imbalanced datasets, while the AUC-ROC plot (Figure~\ref{fig:aucroc_vs_clusters}) highlights stability with scores around 97-99\% across cluster variations, offering valuable insights for optimizing anomaly detection in real-world IoT applications.

% \begin{figure}[!ht]
% \centering
% \begin{subfigure}{\columnwidth}
% \includegraphics[width=\columnwidth]{pictures/aucpr_vs_clusters.png}
% \caption{AUC-PR}
% \label{fig:aucpr_vs_clusters}
% \end{subfigure}%
% \hfill
% \begin{subfigure}{\columnwidth}
% \includegraphics[width=\columnwidth]{pictures/aucroc_vs_clusters.png}
% \caption{AUC-ROC}
% \label{fig:aucroc_vs_clusters}
% \end{subfigure}
% \caption{AUC-ROC and AUC-PR across four IoT datasets with varying numbers of clusters}
% \label{fig:au_vs_cluster}
% \end{figure}


% Further, Critical Difference (CD) diagrams in Figure~\ref{fig:cd_map} compare XXX with baselines across the datasets, evaluating AUC-PR and AUC-ROC without relying on static tests that have been removed. The diagrams present average scores and ranks computed over the four datasets, with XXX leading in AUC-PR at 99.47\% (rank 1.00), far surpassing ECOD at 83.93\% (rank 16.75) and DeepSVDD at 77.83\% (rank 9.00), while KNN at 99.03\% (rank 3.12) and LOF at 98.74\% (rank 4.50) follow closely. In AUC-ROC, XXX achieves 99.61\% (rank 1.00), outpacing ECOD at 89.16\% (rank 16.75) and DeepSVDD at 77.83\% (rank 9.00), with KNN at 99.13\% (rank 3.25) and LOF at 98.74\% (rank 4.25) trailing, emphasizing its superior detection accuracy due to robust feature selection and transformation.





% \begin{figure}[!ht]
% \centering
% \begin{subfigure}{\columnwidth}
% \includegraphics[width=\columnwidth]{pictures/CDDiagram_AUCPR.png}
% \caption{AUCPR}
% \label{fig:cd_aucpr}
% \end{subfigure}%
% \hfill
% \begin{subfigure}{\columnwidth}
% \includegraphics[width=\columnwidth]{pictures/CDDiagram_AUCROC.png}
% \caption{AUCROC}
% \label{fig:cd_aucroc}
% \end{subfigure}
% \caption{AUC-ROC and AUC-PR performance comparison across four IoT datasets using Critical Difference diagrams}
% \label{fig:cd_map}
% \end{figure}


% \newpage
% \subsubsection{Detection Performance on Different Anomaly Types and Noise Levels}
% \label{ssubsec:RQ2}

% This section evaluates the detection performance of the proposed XXX algorithm across various anomaly types and noise levels, comparing it against several baseline methods. The performance is quantified using the Area Under the Receiver Operating Characteristic Curve (AUC-ROC), a widely accepted metric for assessing the ability of a model to distinguish between normal and anomalous instances. The results are visualized in Figure~\ref{fig:noise_vs_tp}, which comprises two subfigures: Figure~\ref{fig:noise_levels} for noise level analysis and Figure~\ref{fig:type_levels} for anomaly type analysis.

% Noise Level and Anomaly Type Analysis: Figure~\ref{fig:noise_levels} illustrates the average AUC-ROC performance under different noise levels (1\%, 3\%, and 5\%), representing the proportion of corrupted data that tests the algorithms' robustness. XXX (AVG-Rank: 1.00) leads with scores around 90\% at 1\% noise, dropping slightly to 85\% at 5\% noise, showcasing exceptional stability. KNN (AVG-Rank: 2.83) follows with a strong 88\% at 1\% noise, declining to 83\% at 5\% noise, indicating good adaptability with a noticeable drop. In contrast, DeepSVDD (AVG-Rank: 8.17) lags with scores around 75\% across noise levels, reflecting lower robustness.  \\
% Figure~\ref{fig:type_levels} compares performance across Global Type, Cluster Type, and Local Type anomaly patterns. XXX (AVG-Rank: 1.50) excels with near 100\% for Global and Cluster Types and 92.5\% for Local Type, demonstrating versatility. LOF (AVG-Rank: 6.58) achieves around 100\% for Global and Cluster Types but drops to 90\% for Local Type, while KNN (AVG-Rank: 7.58) reaches 100\% for Global Type but falls to 82.5\% for Local Type, showing weakness in localized anomalies.\\
% Detailed Insights: The bar plots collectively demonstrate that XXX outperforms the baseline methods across both noise levels and anomaly types, as reflected in its lowest average ranks (1.00 for noise levels and 1.50 for anomaly types). This superiority is attributed to its robust feature selection and transformation process, which effectively mitigates the impact of noise and adapts to various anomaly patterns. 

% \begin{figure}[!ht]
% \centering
% \begin{subfigure}{\columnwidth}    
%     \includegraphics[width=\columnwidth]{pictures/noise_levels.png} 
%     \caption{Noise Levels: Average AUC-ROC performance across different noise percentages for XXX and top baseline methods.}
%     \label{fig:noise_levels}
% \end{subfigure}%
% \hfill
% \begin{subfigure}{\columnwidth}
%     \includegraphics[width=\columnwidth]{pictures/type_levels.png} 
%     \caption{Anomaly Types: Average AUC-ROC performance across Global Type, Cluster Type, and Local Type for XXX and top baseline methods.}
%     \label{fig:type_levels} 
% \end{subfigure}
% \caption{AUC-ROC performance comparison of XXX and the top six baseline methods across different noise levels and anomaly types.}
% \label{fig:noise_vs_tp}
% \end{figure}

% For further details on rankings, parameter configurations, and additional experimental results, please see Section~\ref{exp:more_detail}.

% \subsubsection{Runtime Efficiency and Scalability (RQ3)}
% \label{ssubsec:RQ3}

% The experimental results highlight the exceptional runtime performance of the XXX algorithm, which significantly outperforms the average runtime of other baseline methods in both training and testing phases. This efficiency is further supported by its scalability, as demonstrated in Table~\ref{tab:scale_data}, which presents the runtime performance (in seconds) of the algorithm across varying numbers of samples (\( N \)) and dimensions (\( d \)). The algorithm maintains efficient execution times even as the dataset size and dimensionality increase substantially. However, limitations are observed when handling extremely large datasets, such as \( d = 10,000 \) and \( N = 1,000,000 \), where an \textbf{Out of Memory (OoM)} error occurs. This is attributed to the memory constraints of the evaluation machine, equipped with only 64 GB of RAM, while the dataset demands over 40 GB of memory.

% \begin{table}[!ht]
% \caption{Runtime performance (in seconds) of XXX for different values of \( N \) and \( d \).}
% \centering
% \resizebox{\linewidth}{!}{%
% \begin{tabular}{|c|c|c|c|c|}
% \hline
% \( N \backslash d \) & \( 10 \) & \( 10^2 \) & \( 10^3 \) & \( 10^4 \) \\
% \hline
% \( 10^3 \)   & \( 1.3 \times 10^{-3} \) & \( 1.367 \times 10^{-1} \) & \( 1.1709 \) & \( 6.4533 \times 10^2 \) \\
% \( 10^4 \)   & \( 9.4 \times 10^{-3} \) & \( 2.329 \times 10^{-1} \) & \( 3.3218 \) & \( 7.3267 \times 10^2 \) \\
% \( 10^5 \)   & \( 3.35 \times 10^{-2} \) & \( 1.4780 \) & \( 1.9217 \times 10^1 \) & \( 2.9591 \times 10^3 \) \\
% \( 10^6 \)   & \( 4.915 \times 10^{-1} \) & \( 1.8848 \times 10^1 \) & \( 2.3232 \times 10^2 \) & OoM \\
% \hline
% \end{tabular}%
% }
% \label{tab:scale_data}
% \end{table}

% FFigure~\ref{fig:radar_chart} provides a visual comparison of the performance metrics between the proposed model (XXX) and the baseline methods, focusing on standardized training and testing ranks, as well as train and test times, with the radar chart illustrating key insights: the blue line representing XXX shows a rank of 1.50 for Train Rank and 1.50 for Test Rank, significantly lower than the median baseline performance (orange shaded area with ranks around 6-7), indicating superior performance in both training and testing phases; additionally, the Train Time and Test Time for XXX fall within the inner shaded region, suggesting faster training and testing times compared to the median baseline, which extends further outward, collectively highlighting XXX's efficiency and accuracy advantages. The radar chart uses a normalized scale from 0 to 1, where values closer to the center indicate better performance (lower rank or shorter time), and the blue line (XXX) consistently remains inside the orange shaded area (Median of Baselines), underscoring XXX's overall advantage in terms of both efficiency and effectiveness across all evaluated metrics, reinforcing the algorithm's capability to handle large-scale data efficiently while maintaining high accuracy.

% \begin{figure}[!ht]
% \centering  
% \includegraphics[width=\columnwidth]{pictures/radar_chart.png} 
% \caption{Comparison of standardized training and testing ranks, and train/test times between XXX and baseline methods. The blue line represents XXX, while the orange shaded area indicates the median performance of baselines.}
% \label{fig:radar_chart}
% \end{figure}











% \subsection{Competior}
% To comprehensively assess the effectiveness of our proposed method, \textbf{XXX}, we compare it against several recent state-of-the-art approaches published from 2022 onward. These methods span a range of machine learning and deep learning techniques applied to IoT intrusion detection, including optimization-based models~\cite{aktar2024robust}, genetic algorithms~\cite{saiyed2024genetic}, ensemble frameworks~\cite{golestani2024device}, and unsupervised learning techniques~\cite{segurola2024unsupervised}. The evaluation is conducted on four widely-used public datasets: \textbf{N-BaIoT}, \textbf{ToN-IoT}, \textbf{CICIoT2023}, and \textbf{BoT-IoT}.


% Table~\ref{tab:ids_comparison} reports the performance of all methods in terms of standard classification metrics: \textbf{Accuracy (ACC)}, \textbf{Precision (PRE)}, \textbf{Recall (REC)}, and \textbf{F1-score (F1)}. Our method, XXX, achieves consistently competitive and often superior results across all datasets and metrics.

% On the \textbf{N-BaIoT dataset}, XXX achieves a perfect score of 100\% across all four metrics. While other recent works~\cite{golestani2024device,10839389} report competitive F1-scores (e.g., 92--99\%), none surpass XXX in balanced and complete performance.

% On the \textbf{ToN-IoT dataset}, XXX maintains high performance with 94\% accuracy and an F1-score of 96\%, outperforming other approaches that show broader performance variation~\cite{aktar2024robust, karthikeyan2024firefly, saiyed2024genetic}. Notably, although~\cite{saiyed2024genetic} reports similar accuracy, XXX achieves better consistency across precision and recall.

% In the more challenging \textbf{CICIoT2023 dataset}, which contains complex and imbalanced traffic patterns, XXX delivers 100\% precision and a solid 92\% F1-score. While some prior works such as~\cite{parfenov2024investigation, phan2024comparison} reach comparable accuracy ranges, they lack consistency across precision and recall, and often do not report complete metric sets.

% On the \textbf{BoT-IoT dataset}, XXX reaches 95\% accuracy with 99\% precision and a 97\% F1-score, closely rivaling the strongest reported results~\cite{koroniotis2019towards, golestani2024exploring}. Importantly, XXX achieves stable performance across all evaluation criteria, demonstrating robustness in the presence of noisy and imbalanced IoT data.

% In summary, the results validate the \textbf{superior detection capability and generalizability} of XXX across different datasets and attack types. Its strength lies in effectively preserving local data structure, mitigating the influence of outliers, and maintaining high recall—key requirements for practical anomaly detection systems in dynamic IoT environments.
% \begin{table*}[!ht]
% \caption{Performance comparison of our proposed method (XXX) and recent research on four IoT datasets using Accuracy, Precision, Recall, and F1-score.}
% \centering
% \small % Smaller font for better scaling
% \setlength{\tabcolsep}{6pt} % Increase column spacing
% \resizebox{\textwidth}{!}{%
% \begin{tabular}{|c|c|c|c|c|c|c|}
% \hline
% \textbf{Dataset} & \textbf{Method} & \textbf{ACC (\%)} & \textbf{PRE (\%)} & \textbf{REC (\%)} & \textbf{F1 (\%)} \\ 
% \hline
% {\textbf{N-BaIoT}} & \textbf{XXX
% (Our model)} & 100.00 & 100.00 & 100.00 & 100.00 \\ 
% & Ref \cite{golestani2024device} & 88.00--99.00 & \text{--} & \text{--} & 92.00--99.00 \\
% & Ref \cite{10839389} & \text{--} & 95.00 & 91.00 & 88.00 \\
% \hline
% {\textbf{ToN-IoT}} & \textbf{XXX} & 94.00 & 98.00 & 95.00 & 96.00 \\ 
% & Ref \cite{aktar2024robust} & 88.00--99.00 & 84.00--98.00 & 72.00--99.00 & 76.00--99.00 \\
% & Ref \cite{karthikeyan2024firefly} & 85.00--94.00 & \text{--} & \text{--} & \text{--} \\
% & Ref \cite{saiyed2024genetic} & 95.00 & 95.00 & 95.00 & 94.00 \\
% \hline
% {\textbf{CICIoT2023}} & \textbf{XXX} & 85.00 & 100.00 & 85.00 & 92.00 \\
% & Ref \cite{parfenov2024investigation} & 95.00 & \text{--} & \text{--} & \text{--} \\
% & Ref \cite{golestani2024device} & 84.00--90.00 & \text{--} & \text{--} & 88.00--92.00 \\
% & Ref \cite{phan2024comparison} & 83.00--92.00 & 82.00--91.00 & 82.00--93.00 & \text{--} \\
% \hline
% {\textbf{BoTIoT}} & \textbf{XXX} & 95.00 & 99.00 & 95.00 & 97.00 \\
% & Ref \cite{koroniotis2019towards} & 97.00 & 100.00 & 97.00 & \text{--} \\
% & Ref \cite{segurola2024unsupervised} & \text{--} & \text{--} & 97.00 & 98.00 \\
% & Ref \cite{9993869} & \text{--} & \text{--} & 96.00 & 96.00 \\
% & Ref \cite{golestani2024exploring} & 45.00--97.00 & 100.00 & 45.00--97.00 & 62.00--98.00 \\
% \hline
% \end{tabular}
% }
% \label{tab:ids_comparison}
% \end{table*}

% % ========================
% \section{Conclusion}
% % ========================
% This study introduces a novel algorithm called Local One-Class Null Foley-Sammon Transformation (XXX), designed to address the one-class novelty detection challenge within the realm of network intrusion detection for IoT environments. The XXX algorithm enhances the traditional Null Foley-Sammon Transformation (NFST) by applying K-means clustering to simulate pseudo-classes and leveraging Singular Value Decomposition (SVD) for stable subspace projection. This approach preserves both local and global data structures, mitigating the impact of singular points and improving accuracy when dealing with complex IoT network traffic. Moreover, XXX employs a robust projection framework to ensure accuracy and utilizes efficient matrix operations to reduce computational costs. Extensive experiments on various public datasets, including CICIoT2023, N-BaIoT, BoTIoT, and ToN-IoT, demonstrate the superiority of XXX compared to other baseline algorithms across diverse attack scenarios, achieving AUC-ROC scores of 99.61\% and AUC-PR scores of 99.47\%.

% Besides representing a significant advancement in addressing the one-class novelty detection challenge, this research presents several open research directions. One key direction is the investigation of parallelization and hardware acceleration to enhance the computational efficiency of XXX, particularly in clustering and SVD computations. Another critical area for future work is detecting novel patterns in benign/normal traffic. Developing labeled datasets for benign traffic and refining methods to identify novel benign patterns will strengthen IDS performance in dynamic IoT network environments. 
% Additionally, future research will focus on developing an adaptive mechanism for optimal threshold selection in novelty detection. By allowing thresholds to adjust dynamically based on dataset characteristics, it can further improve detection accuracy and minimize false positives, ensuring the model can effectively differentiate between novel and known traffic patterns.

% % ========================


% \bibliographystyle{elsarticle-num}
% \bibliography{references}


% \newpage
% \appendix
% \onecolumn


% \section{Detailed Supplementary Theory}
% \label{ap1:supp_theory}

% This section provides additional theoretical insights to clarify the technical and methodological aspects related to the presented experiments. It focuses on explaining core concepts, model approaches, and characteristics of the IoT datasets used.

% \subsection{Datasets}
% \label{ap1:datasets}

%  In this study, we thoroughly evaluate our proposed method
%  on four widely used datasets in the attack detection field: BoT
% IoT, CIC-IoT2023, N-BaIoT, and ToN-IoT. These datasets are
%  specifically designed to simulate IoT network environments and
%  include a range of up-to-date attacks. A summary of each
%  dataset is provided in Table 3, and further details are described
%  below:
% \begin{table}[h]
% \centering
% \caption{Summary of IoT Datasets Used for Evaluation}
% \label{tab:datasets}
% \begin{tabular}{|l|c|c|c|}
% \hline
% \textbf{Dataset}   & \textbf{Classes} & \textbf{Features} \\
% \hline
% CICIoT2023         & 6                & 46                       \\
% N-BaIoT            & 9                & 115                  \\
% BoTIoT             & 6                & 21                \\
% ToNIoT             & 10               & 21                    \\
% \hline
% \end{tabular}
% \end{table}

 
% \textbf{CICIoT2023} : This Dataset developed by the Canadian Institute for Cybersecurity, CICIoT2023 is a comprehensive dataset designed for IoT intrusion detection research~\cite{neto2023ciciot2023}. It includes traffic from various smart devices such as thermostats and cameras under both normal operations and several attack scenarios (e.g., DoS, DDoS, botnets, and reconnaissance). The dataset offers 46 network flow features, including packet sizes, protocol types, and durations. With a large sample size and diverse attack types, CICIoT2023 enables robust assessment of OC-NFST in detecting various high-dimensional and noisy anomalies.



% \textbf{N-BaIoT} : The dataset targets anomaly detection in smart home environments~\cite{8490192}. It contains traffic from connected devices such as cameras and sensors under both benign and malicious conditions. Comprising 115 features—including packet counts, byte sizes, and connection timings—the dataset is particularly useful for testing clustering-based and noise-tolerant methods. Noise injection simulates real-world imperfections, testing OC-NFST’s effectiveness in identifying both global and local anomalies.


% \textbf{BoTIoT}: This datasets is specifically curated for evaluating botnet detection in IoT ecosystems~\cite{koroniotis2019towards}. It includes time-series data from smart appliances and connected vehicles under normal usage and botnet-induced attacks such as scanning and DDoS. Featuring 21 attributes like packet lengths and inter-arrival times, the dataset focuses on time-dependent anomalies, making it ideal for evaluating OC-NFST’s performance on dependency-based detection in dynamic IoT scenarios.



% \textbf{ToNIoT}~\cite{fesz-dm97-19}: This dataset was provided by the Cyber Range Lab of UNSW Canberra to support the development of artificial intelligence-based security solutions for industrial IoT networks. It consists of benign network traffic generated by legitimate IoT devices under normal operating conditions. Moreover, the dataset includes diverse IoT applications and communication protocols commonly used in IoT deployments.

% \vspace{0.5em}


% \subsection{Additional Demonstration of Synthetic Anomalies}
% \label{ap1:ano_types}

% To evaluate the performance of our approach, we test it on 47 core benchmark datasets, which is similar to the datasets used in the original ADBench paper \cite{han2022adbench}. This allows us to assess how well the proposed methods work across a wide range of real-world scenarios and anomaly types. Below, we describe the generation process for four common types of anomalies.

% \textbf{Definition and Generation Process of Four Common Types of Anomalies} \cite{han2022adbench}:

% \begin{itemize}
%     \item \textbf{Local Anomalies:} Local anomalies are points that deviate significantly from their local neighborhoods. To generate local anomalies, we first create synthetic normal samples using the GMM \cite{GMM}. Then, we scale the covariance matrix by a scaling parameter $\alpha = 5$ to generate local anomalies, which alters the covariance structure of the normal data.
    
%     \item \textbf{Global Anomalies:}  Global anomalies are data points that are significantly different from the normal data. These anomalies are generated using a uniform distribution $\text{Unif}(min(X_k), max(X_k))$, where the boundaries are defined by the minimum and maximum values of an input feature $X_k$. The scaling parameter $\alpha = 1.1$ controls how extreme these anomalies are.
    
%     \item \textbf{Dependency Anomalies:} Dependency anomalies refer to samples that do not follow the dependency structure of the normal data. In this case, the input features of the anomalies are assumed to be independent of each other. We apply the Vine Copula method  \cite{VinCP} to model the dependency structure of the original normal data. The probability density function (PDF) for generating anomalies is set to complete independence by removing the modeled dependency. Notably, because the Vine Copula \cite{VinCP} method can be computationally expensive when working with large datasets, we only use 2,000 normal samples to fit this model. After fitting the model, we use Kernel Density Estimation (KDE) \cite{KDE} to estimate the PDF of the features and generate independent anomaly samples.

%     \item \textbf{Cluster Anomalies:} Cluster anomalies, also known as group anomalies, exhibit similar characteristics to clusters. To generate cluster anomalies, we adjust the mean feature vector of the normal samples by a scaling parameter $\alpha = 5$, i.e., shifting the mean of the GMM \cite{GMM}. This parameter controls the distance between anomaly clusters and the normal data. Then, the anomalies are generated by applying the scaled GMM \cite{GMM}.
% \end{itemize}



% Building on the approach from ADBench, we generate "realistic" synthetic data using a generative model like Gaussian Mixture Model (GMM) \cite{GMM}, trained on the normal samples of a benchmark dataset. The original anomalies are discarded since their exact types are unknown. By adjusting the parameters of the generative model, we can create both normal samples and different types of anomalies. This method adds complexity to the generated anomalies, making the synthetic data better reflect the challenges of real-world anomaly detection tasks.

% \subsection{Baseline Algorithms}
% \label{ap1:ap_baselines}

% In this section, we present the baseline algorithms employed for anomaly detection. These algorithms are categorized into two groups: \textit{Shallow Machine Learning-based} and \textit{Deep Learning-based} outlier detection algorithms. 

% \subsubsection{Machine Learning-based Outlier Detection Algorithms}

% \begin{itemize}
%     \item \textbf{Clustering-Based Local Outlier Factor (CBLOF)} \cite{CBLOF}: CBLOF identifies anomalies by first partitioning the data into clusters and then evaluating the distance between the data points and the cluster centroids. Instances that exhibit significant distance from their respective clusters are considered outliers.

%     \item \textbf{K-Nearest Neighbors (KNN)} \cite{KNN}: KNN calculates anomaly scores based on the distance to the $k$-th nearest neighbor. Data points with larger distances from their nearest neighbors are assigned higher anomaly scores, suggesting that they are potential outliers.

%     \item \textbf{Isolation Forest (IForest)} \cite{IForest}: Isolation Forest isolates observations by recursively partitioning the data using random features and split values. The anomalies, which are few and distinct, are isolated more efficiently, resulting in shorter path lengths.

%     \item \textbf{One-Class Support Vector Machine (OCSVM)} \cite{OCSVM}: OCSVM constructs a decision boundary in a high-dimensional space, separating normal data points from outliers by maximizing the margin between the origin and the normal data points. The decision boundary thus defines the region where anomalies are unlikely to reside.

%     \item \textbf{Local Outlier Factor (LOF)} \cite{LOF}: LOF measures the local density deviation of a data point in relation to its neighbors. Points exhibiting substantially lower density compared to their neighbors are considered anomalies.

%     \item \textbf{Histogram-Based Outlier Detection (HBOS)} \cite{HBOS}: HBOS assumes the independence of the features and constructs histograms for each feature. Anomalies are identified based on the deviation of a data point from the expected distribution represented by these histograms.

%     \item \textbf{Principal Component Analysis (PCA)} \cite{PCA}: PCA is employed to reduce the dimensionality of the data by projecting it onto a lower-dimensional space. Anomalies are detected by evaluating the reconstruction error, with larger errors indicating a poor fit to the reduced space and suggesting outlier behavior.

%     \item \textbf{Subspace Outlier Detection (SOD)} \cite{SOD}: SOD detects outliers by examining various subspaces of the feature space. Anomalies are identified based on their deviation from expected behavior in one or more of these subspaces.

%     \item \textbf{Connectivity-Based Outlier Factor (COF)} \cite{COF}: COF measures the chaining distance of a data point and its k-th nearest neighbors. Data points with unusually large or small chaining distances relative to their neighbors are classified as anomalies.

%     \item \textbf{Lightweight Online Detector of Anomalies (LODA)} \cite{LODA}: LODA is an ensemble-based anomaly detection algorithm particularly suitable for real-time anomaly detection in streaming data. It is capable of adapting to concept drift and large-scale data environments.

%     \item \textbf{Empirical Cumulative Distribution-Based Outlier Detection (ECOD)} \cite{ECOD}: ECOD uses the empirical cumulative distribution function (ECDF) to model the distribution of each feature. Anomalies are detected by identifying data points that lie in the extreme tails of the distribution.

%     \item \textbf{Copula-Based Outlier Detector (COPOD)} \cite{COPOD}: COPOD is a model-free, highly interpretable anomaly detection technique that uses copula-based models to capture the dependence structure between features. Detects anomalies in multivariate data without requiring hyperparameter tuning.
% \end{itemize}

% \subsubsection{Deep Learning-based Outlier Detection Algorithms}

% \begin{itemize}
%     \item \textbf{Deep Support Vector Data Description (DeepSVDD)} \cite{DeepSVDD}: DeepSVDD uses a deep neural network to learn representations of the data while minimizing the volume of a hypersphere that encloses these representations. Data points that lie far from the center of this hypersphere are identified as anomalies.

%     \item \textbf{AutoEncoder} \cite{Autoencoder}: AutoEncoders are neural networks designed to learn efficient representations of data through a process of encoding and decoding. Anomalies are detected based on reconstruction errors, with higher errors indicating instances that do not conform to the learned data distribution.

%     \item \textbf{Deviation Networks (DevNet)} \cite{DevNet}: DevNet is a deep learning model that enforces a deviation-based loss function to penalize large deviations in the learned representations. This method is particularly effective in identifying outliers by ensuring that anomalous instances exhibit substantial deviations from the learned data distribution.

%     \item \textbf{Learning Unsupervised Anomalies with Neural Representation learning (LUNAR)} \cite{LUNAR}: LUNAR is a deep learning-based, neural network approach designed for unsupervised anomaly detection. It learns a data distribution through representation learning and detects anomalies by identifying instances that deviate significantly from the normal patterns learned. This method is particularly effective in high-dimensional, complex data environments where labeled data is unavailable.



%     \item \textbf{Autoencoder-based One-Class Support Vector Machine (AE1SVM)} \cite{AE1SVM}: AE1SVM is a hybrid method combining AutoEncoder-based feature learning with Support Vector Machines for enhanced anomaly detection. This approach leverages the reconstruction errors from the AutoEncoder in conjunction with SVM’s classification capabilities to detect outliers.

%     \item \textbf{Adversarial Learned Anomaly Detection (ALAD)} \cite{ALAD}: ALAD (additional details to be included) is a deep learning model that utilizes adversarial learning techniques for anomaly detection. The architecture is designed to capture the subtle differences between normal and anomalous instances by learning the underlying data distribution.
% \end{itemize}

% \subsection{Evaluation Metrics}
% \label{ap1:EM}
% To comprehensively assess the performance of the proposed method, we adopt a suite of evaluation metrics that capture various aspects of binary classification performance, especially relevant to anomaly detection in imbalanced and noisy IoT network traffic. The formulas and interpretations are as follows:

% \begin{itemize}
%     \item \textbf{Accuracy (ACC):} Measures the proportion of correctly predicted instances among all samples.
%     \[
%     \text{ACC} = \frac{\text{TP} + \text{TN}}{\text{TP} + \text{TN} + \text{FP} + \text{FN}}
%     \]

%     \item \textbf{True Positive Rate (TPR) or Recall:} Indicates the ability to correctly identify actual anomalies.
%     \[
%     \text{TPR} = \frac{\text{TP}}{\text{TP} + \text{FN}}
%     \]

%     \item \textbf{False Positive Rate (FPR):} Represents the proportion of normal instances incorrectly classified as anomalies.
%     \[
%     \text{FPR} = \frac{\text{FP}}{\text{FP} + \text{TN}}
%     \]

%     \item \textbf{Positive Predictive Value (PPV) or Precision:} Measures how many of the instances predicted as anomalies are actually anomalous.
%     \[
%     \text{PPV} = \frac{\text{TP}}{\text{TP} + \text{FP}}
%     \]

%     \item \textbf{F1-score:} Harmonic mean of precision and recall, providing a balanced metric for imbalanced datasets.
%     \[
%     \text{F1-score} = \frac{2 \cdot \text{TPR} \cdot \text{PPV}}{\text{TPR} + \text{PPV}}
%     \]

%     \item \textbf{Matthews Correlation Coefficient (MCC):} A correlation-based metric that considers all four values in the confusion matrix and is robust to imbalanced datasets.
%     \[
%     \text{MCC} = \frac{\text{TP} \cdot \text{TN} - \text{FP} \cdot \text{FN}}{\sqrt{(\text{TP} + \text{FP})(\text{TP} + \text{FN})(\text{TN} + \text{FP})(\text{TN} + \text{FN})}}
%     \]

%     \item \textbf{Area Under the ROC Curve (AUC-ROC):} Quantifies the model’s ability to distinguish between classes across various thresholds. It is computed as the area under the ROC curve, which plots TPR versus FPR.
%     \[
%     \text{AUC-ROC} = \sum_{i=1}^{n-1} \frac{\text{TPR}_{i+1} + \text{TPR}_i}{2} \cdot (\text{FPR}_{i+1} - \text{FPR}_i)
%     \]
%     A value closer to 1 indicates stronger class separation capability.

%     \item \textbf{Area Under the Precision-Recall Curve (AUC-PR):} Captures the trade-off between precision and recall, especially informative in scenarios with significant class imbalance. It is typically approximated using the trapezoidal rule:
%     \[
%     \text{AUC-PR} = \sum_{i=1}^{n-1} (\text{Recall}_{i+1} - \text{Recall}_i) \cdot \frac{\text{Precision}_{i+1} + \text{Precision}_i}{2}
%     \]
%     Higher AUC-PR values reflect better performance in identifying positive (anomalous) instances.
% \end{itemize}

% This comprehensive set of metrics provides a robust foundation for evaluating the accuracy, reliability, and robustness of our proposed method under diverse and challenging IoT traffic scenarios.


% \section{Experiments with More Details and Results}
% \label{ap2:experiments}

% This section details the experimental setup and results, analyzing model performance across IoT datasets under varying noise levels and anomaly types. Accuracy and AUCROC scores, reported in Tables \ref{tab:model_noise_perf} and \ref{tab:aucroc_results}, highlight OurModel’s superiority, achieving up to 96.29\% accuracy on BIoT (1\% noise) and 100\% AUCROC for global anomalies. A survey of recent IDS research in Table \ref{tab:ids_overview} compares 11 studies, showing diverse approaches with accuracies from 45\% to 99\%, contextualizing our method’s advancements.


% \subsection{Expriments results with more details}
% \label{ap2:more_detail}
% This section presents a detailed evaluation of various models across multiple IoT datasets under different noise levels and anomaly types. Table \ref{tab:model_noise_perf} reports the model performance under varying noise conditions (1\%, 3\%, 5\%) on BIoT, CICIoT2023, NBaIoT, and ToNIoT datasets, highlighting the robustness of each model. Table \ref{tab:aucroc_results} provides AUCROC scores across local, global, and cluster anomaly types, showcasing the models' effectiveness in detecting different anomaly patterns.


% Table \ref{tab:model_noise_perf} summarizes the performance of various models under different noise levels (1\%, 3\%, 5\%) across four IoT datasets: BIoT, CICIoT2023, NBaIoT, and ToNIoT. OurModel consistently outperforms others, achieving the highest scores across all noise levels and datasets, with top accuracies like 96.29\% (BIoT, 1\% noise) and 89.53\% (ToNIoT, 1\% noise). KNN follows closely, particularly strong in BIoT and ToNIoT, while AutoEncoder and CBLOF show competitive results in specific scenarios. Other models, such as DevNet and COPOD, exhibit lower performance, especially under higher noise levels. 

% \begin{table*}[!ht]
% \caption{Model performance across different noise levels (1\%, 3\%, 5\%) on various IoT datasets.}
% \centering
% \resizebox{\textwidth}{!}{%
% \begin{tabular}{|c|ccc|ccc|ccc|ccc|}
% \hline
% \textbf{Model} & \multicolumn{3}{c|}{\textbf{BIoT}} & \multicolumn{3}{c|}{\textbf{CICIoT2023}} & \multicolumn{3}{c|}{\textbf{NBaIoT}} & \multicolumn{3}{c|}{\textbf{ToNIoT}} \\
% \cline{2-13}
%  & 1\% & 3\% & 5\% & 1\% & 3\% & 5\% & 1\% & 3\% & 5\% & 1\% & 3\% & 5\% \\
% \hline
% OurModel & \textcolor{red}{96.29} & \textcolor{red}{93.77} & \textcolor{red}{92.90} & \textcolor{red}{95.37} & \textcolor{red}{93.93} & \textcolor{red}{93.35} & \textcolor{red}{79.24} & \textcolor{red}{77.36} & \textcolor{red}{75.35} & \textcolor{red}{89.53} & \textcolor{red}{83.92} & \textcolor{red}{81.47} \\
% KNN & \textcolor{blue}{96.00} & \textcolor{blue}{90.63} & \textcolor{blue}{88.24} & \textcolor{blue}{93.11} & \textcolor{blue}{92.43} & 90.09 & 78.15 & \textcolor{blue}{76.02} & \textcolor{blue}{74.81} & \textcolor{blue}{86.31} & \textcolor{blue}{83.54} & \textcolor{blue}{79.79} \\
% AutoEncoder & 95.66 & 89.62 & 86.45 & 92.46 & 91.64 & \textcolor{blue}{ 91.24}& 55.89 & 21.90 & 66.72 & 66.32 & 67.32 & 71.73 \\
% CBLOF & 87.06 & 86.50 & 85.99 & 90.68 & 90.65 & 89.90 & 75.28 & 75.88 & 74.70 & 74.35 & 73.92 & 76.69 \\
% AE1SVM & 92.43 & 91.41 & 86.94 & 90.19 & 90.77 & 90.44 & 56.06 & 64.31 & 48.93 & 54.71 & 56.57 & 52.80 \\
% DevNet & 64.14 & 63.63 & 64.11 & 85.86 & 85.82 & 85.85 & \textcolor{blue}{ 78.28} & 64.49 & 64.49 & 76.31 & 76.36 & 76.30 \\
% PCA & 90.19 & 88.21 & 86.07 & 89.97 & 89.44 & 88.62 & 69.16 & 66.30 & 61.66 & 53.64 & 52.46 & 52.10 \\
% IForest & 84.24 & 80.56 & 82.09 & 92.29 & 92.18 & 91.66 & 30.87 & 29.73 & 27.18 & 76.57 & 75.38 & 75.80 \\
% LODA & 76.00 & 91.55 & 79.47 & 83.50 & 81.36 & 86.99 & 47.32 & 41.75 & 63.17 & 56.56 & 69.65 & 62.53 \\
% DeepSVDD & 87.27 & 93.63 & 79.28 & 86.40 & 90.33 & 89.47 & 62.17 & 61.35 & 57.68 & 63.10 & 62.24 & 63.24 \\
% HBOS & 85.71 & 84.57 & 83.79 & 94.13 & 80.96 & 78.59 & 31.01 & 29.72 & 28.84 & 82.05 & 81.49 & 80.74 \\
% LOF & 92.38 & 65.33 & 61.21 & 88.92 & 82.26 & 75.06 & 76.46 & 57.45 & 56.32 & 82.68 & 76.36 & 70.72 \\
% MO\_GAAL & 39.89 & 44.38 & 45.17 & 54.96 & 55.10 & 55.68 & 59.74 & 31.31 & 35.14 & 62.32 & 77.88 & 77.78 \\
% OCSVM & 85.75 & 83.69 & 81.42 & 84.39 & 84.01 & 83.40 & 77.63 & 76.66 & 30.47 & 48.38 & 48.11 & 48.31 \\
% ALAD & 53.62 & 53.10 & 49.96 & 50.57 & 51.17 & 51.14 & 43.33 & 43.30 & 45.49 & 51.90 & 51.52 & 52.41 \\
% COPOD & 40.73 & 40.66 & 40.63 & 44.90 & 44.88 & 44.86 & 31.88 & 31.77 & 31.61 & 29.81 & 29.79 & 29.75 \\
% ECOD & 42.34 & 42.27 & 42.24 & 34.86 & 34.86 & 34.84 & 30.44 & 30.25 & 30.11 & 23.59 & 23.57 & 23.55 \\
% \hline
% \end{tabular}
% }
% \label{tab:model_noise_perf}
% \end{table*}

% Table \ref{tab:aucroc_results} presents AUCROC scores for models across local, global, and cluster anomaly types on the same IoT datasets. OurModel again leads, achieving near-perfect scores (e.g., 100\% for global and cluster anomalies in BIoT, NBaIoT, and ToNIoT). KNN, AutoEncoder, and PCA show strong performance, particularly for global and cluster anomalies, with scores often exceeding 99\%. Models like DevNet and ALAD struggle with local anomalies, while IForest and HBOS perform well in specific cases, such as global anomalies in BIoT and NBaIoT.


% \begin{table}[!ht]
% \caption{AUCROC Scores for Different Models across Anomaly Types and Datasets}
% \centering
% \resizebox{\textwidth}{!}{%
% \begin{tabular}{|c|ccc|ccc|ccc|ccc|}
% \hline
% \textbf{Model} & \multicolumn{3}{c|}{\textbf{BoT-IoT}} & \multicolumn{3}{c|}{\textbf{CICIoT2023}} & \multicolumn{3}{c|}{\textbf{N-BaIoT}} & \multicolumn{3}{c|}{\textbf{ToN-IoT}} \\
%  & Local & Global & Cluster & Local & Global & Cluster & Local & Global & Cluster & Local & Global & Cluster \\
% \hline
% OurModel & \textcolor{red}{95.32} & \textcolor{red}{100} & \textcolor{red}{100} & 91.29 & \textcolor{red}{99.94} & \textcolor{red}{100} & \textcolor{red}{92.76} & \textcolor{red}{100} & \textcolor{red}{99.99} & \textcolor{red}{93.16} &\textcolor{red}{100} & \textcolor{red}{ 100} \\
% KNN & 90.11 & 99.93 & 99.93 & 70.90 & 99.76 & 99.73 & 83.62 & 99.59 & \textcolor{blue}{99.85} & 85.89 & 99.93 & 99.95 \\
% AutoEncoder & 89.58 & 99.83 & 99.86 & 95.44 & 99.35 & 99.56 & 89.40 & 99.92 & \textcolor{blue}{99.85} & \textcolor{blue}{91.11} & 99.89 & 99.94 \\
% CBLOF & 89.49 & 99.72 & 99.93 & 69.51 & 99.69 & 99.63 & 85.07 &  \textcolor{blue}{ 99.99} & 99.70 & 85.38 & 99.93 & \textcolor{blue}{99.96} \\
% AE1SVM & 90.41 & 99.80 & 99.93 & 95.16 & 99.56 &  \textcolor{blue}{99.78} & 89.53 & 99.85 & 99.77 & 89.03 & 99.93 & 99.88 \\
% DevNet & 52.71 & 24.49 & 3.48 & 54.50 & 54.29 & 11.40 & 51.38 & 41.33 & 3.40 & 52.05 & 31.88 & 31.19 \\
% PCA & 87.26 & 99.66 & 99.86 & \textcolor{blue}{95.24} &  \textcolor{blue}{99.78} & \textcolor{blue}{99.78} & 86.34 & \textcolor{red}{100} & 99.92 & 84.93 & 99.88 & 99.88 \\
% IForest & 83.23 & \textcolor{blue}{99.99} & \textcolor{red}{100} & 94.34 & 99.33 & 99.56 & 92.72 & \textcolor{red}{100} & 99.78 & 84.38 & \textcolor{red}{100} & 99.75 \\
% LODA & 89.57 & 99.45 & 99.84 & 82.10 & 96.80 & 99.52 & \textcolor{blue}{92.56} & 99.27 & 99.93 & 76.45 & 99.96 & 98.66 \\
% DeepSVDD & 86.47 & 99.58 & 99.93 & 90.49 & 99.55 & 99.55 & 90.20 & 99.91 & 99.93 & 91.92 & 99.92 & 99.94 \\
% HBOS & 77.33 & 99.93 &   \textcolor{red}{100} & 90.37 & 99.55 & 96.98 & 92.70 & \textcolor{red}{100} & 99.99 & 79.38 & 99.94 & 99.88 \\
% LOF & \textcolor{blue}{92.63} & 99.89 & 99.93 & 90.71 & 99.76 & 99.86 & 88.65 & 99.54 & 99.70 & 90.59 & 99.87 & \textcolor{red}{100} \\
% MO\_GAAL & 55.83 & 99.85 & 99.70 & 65.05 & 66.52 & 99.68 & 62.35 & 67.76 & 99.68 & 65.07 & 57.43 & 97.56 \\
% OCSVM & 90.56 & 99.72 & 99.93 & 89.85 & 99.28 & 99.77 & 92.33 & 99.45 & 99.63 & 88.32 & 99.94 & 99.98 \\
% ALAD & 54.34 & 61.86 & 83.01 & 51.74 & 60.39 & 86.63 & 50.65 & 57.89 & 89.39 & 55.67 & 62.02 & 75.87 \\
% COPOD & 73.89 & 99.76 & 30.07 & 92.45 & 99.73 & 56.28 & 91.25 & \textcolor{red}{100} & 12.68 & 92.76 & \textcolor{blue}{99.99} & 54.17 \\
% ECOD & 78.26 & 99.82 & 48.26 & \textcolor{red}{96.59} & 99.91 & 44.45 & 85.69 & \textcolor{blue}{99.99} & 19.63 & 92.22 & \textcolor{blue}{99.99} & 50.47 \\
% \hline
% \end{tabular}%
% }
% \label{tab:aucroc_results}
% \end{table}


% \subsection{Survey of Competitor Results}
% \label{ap2:competitor_}

% This section provides an overview of recent research on Intrusion Detection Systems (IDSs) for IoT networks, as summarized in Table \ref{tab:ids_overview}. The table compiles key details from 11 studies, including their titles, years, methodologies, datasets, tasks, models, and performance metrics. The surveyed works primarily focus on anomaly detection, binary classification, and novelty detection, utilizing datasets such as BoTIoT, CICIoT2023, N-BaIoT, and TON-IoT. Notable models include DeepSVDD, OCSVM, KANs, and RNN-based approaches, with performance metrics ranging widely (e.g., accuracy from 45\% to 99\%, F1 scores from 62\% to 98\%). These studies highlight diverse approaches to enhancing IoT security, with varying degrees of effectiveness across different datasets and tasks.

% % ===================== Table in full-page =====================
% \onecolumn

% % Trong nội dung:
% \begin{table*}[p] % full-page width
% \centering
% \caption{Overview of recent research on IDSs, summarizing training approach, task, model, datasets, and performance.}
% \label{tab:ids_overview}
% \renewcommand{\arraystretch}{2.1} % tăng khoảng cách dòng
% % \small  % chữ to hơn nếu cần \normalsize hoặc \large
% \scriptsize % Cỡ nhỏ hơn small 
% \setlength{\tabcolsep}{4pt} % giảm khoảng cách giữa cột

% \begin{adjustbox}{width=1\textwidth,center} % có thể tăng lên nếu cần
% \begin{tabular}{@{}c>{\centering\arraybackslash}m{3.5cm}c>{\arraybackslash}m{4.8cm}c>{\centering\arraybackslash}m{2.8cm}c>{\centering\arraybackslash}m{2.5cm}c>{\centering\arraybackslash}m{3.2cm}@{}}
% \toprule
% \textbf{ID} & \textbf{Title} & \textbf{Year} & \textbf{Summary} & \textbf{Dataset} & \textbf{Tasks} & \textbf{Model} & \textbf{Performance} \\ 
% \midrule    

% 1 & Exploring Unsupervised One-Class Classifiers for Lightweight Intrusion Detection in IoT Systems\cite{golestani2024exploring}& 2024 & Compares five one-class classifiers with ML/DL models in IoT & BoTIoT & Anomaly detection & iForest + OCSVM + LOF + DeepSVDD + DROCC & ACC: 45\%-97\%, PRE: 100\%, REC: 45\%-97\%, F1: 62\%-98\% \\

% 2 & Device-Specific Anomaly Detection Models for IoT Systems\cite{golestani2024device} & 2024 & Compare single models and device -specific model & N-BaIoT, CICIoT2023 & Anomaly Detection & DeepSVDD + LOF + RF + DT & ACC: 84\%-99\%, F1: 88\%-99\% \\

% 3 & Classifying IoT Botnet Attacks With Kolmogorov-Arnold Networks: A Comparative Analysis of Architectural Variations \cite{10839389}
%  & 2024 & Kolmogorov-Arnold Networks (KANs) and their variations in detecting IoT botnet attacks & N-BaIoT &  Anomaly Detection & ORIGINAL-KAN & PRE: 95\%, REC: 91\%, F1:88\%\\

% 4 &Robust Anomaly Detection in IoT Networks using Deep SVDD and Contractive Autoencoder\cite{aktar2024robust} & 2024 & Using Deep SVDD and Contractive Autoencoder in Anomaly Detection
%  & TON-IoT &  Anomaly Detection & KMeans+OCSVM+Isolation Forest+DSVDD-CAE & ACC: 88-99\%, PRE: 84-98\%, REC: 72-99\%, F1: 76-99\% \\

% % 5 Chưa xong 
% 5 & Firefy algorithm based WSN‑IoT security enhancement with machine learning for intrusion detection\cite{karthikeyan2024firefly} & 2024 & RNN-based IDS & TON-IoT  & Intrusion Detection OR?  Anomaly Detection  & RNN/LSTM/GRU & ACC: 94–85\% \\ 

% 6 &A Genetic Algorithm and t-Test-Based System for DDoS Attack Detection in IoT Networks\cite{saiyed2024genetic} & 2024 & Genetic Algorithm and t-Test-Based System for DDoS Attack Detection in IoT Networks &TON-IoT &  binary classification & GADAT-ET  & ACC: 95\%, PRE: 95\%, REC: 95\%, F1: 94\% \\

% 7 &Investigation of the impact effectiveness of adversarial data leakage attacks on the machine
% learning models\cite{parfenov2024investigation}  & 2023 & Adversarial attacks in detecting anomalous activity in IoT networks  & CICIoT2023  & binary classification & MLP-Prod + Catboost & ACC: 95\% \\

% 6 &A Genetic Algorithm and t-Test-Based System for DDoS Attack Detection in IoT Networks\cite{saiyed2024genetic} & 2024 & Genetic Algorithm and t-Test-Based System for DDoS Attack Detection in IoT Networks & TON-IoT &  Binary classification & GADAT-ET  & ACC: 95\%, PRE: 95\%, REC: 95\%, F1: 94\% \\

% 7 &Investigation of the impact effectiveness of adversarial data leakage attacks on the machine
% learning models\cite{parfenov2024investigation}  & 2023 & Adversarial attacks in detecting anomalous activity in IoT networks  & CICIoT2023  & binary classification & MLP-Prod + Catboost & ACC: 95\% \\

% 8 & Comparison of Multiple Feature Selection Techniques for Machine Learning-Based Detection of IoT Attacks\cite{phan2024comparison} & 2024 & Experience of feature selection methods & CICIoT2023 & Novelty Detection & K-NN + MLP & ACC:83-92\%, PRE: 82-91\%, REC: 82-93\%\\

% %  9 chưa xong 
% 9 & Towards the Development of Realistic Botnet Dataset in the Internet of Things for Network Forensic Analytics: Bot-IoT Dataset\cite{koroniotis2019towards} & 2019 &   A new dataset & BoTIoT & Binary classification & RNN & ACC: 97\%, PRE: 100\%, REC: 97\% \\

% 10 & Unsupervised Anomaly Detection Approach for Cyberattack Identification\cite{segurola2024unsupervised} & 2024 & Three-step unsupervised anomaly detection (empirical Bayesian scoring, feature selection, model election) & BoTIoT & Anomaly Detection & RSR + EAD  & REC: 97\%, F1:98\%.  \\ 

% 11 & Improving the IDS for BoT-IoT Dataset-Based Machine Learning Classifiers \cite{9993869} & 2022 & Experiences on BoTIoT dataset & BoTIoT & Binary classification  & RF  & REC: 96\%, F1:96\%.  \\ 

% \bottomrule
% \end{tabular}
% \end{adjustbox}
% \end{table*}

% \twocolumn



% % 


% % =============================================================


% % \begin{table*}[!ht]
% % \centering
% % \caption{The overview of recent research on IDSs, summarizing the training approach, task, model, datasets used, and detection performance.}
% % \label{tab:ids_overview}
% % \resizebox{\textwidth}{!}{%
% % \begin{tabular}{|c|p{4.2cm}|c|p{4.5cm}|c|c|p{3.5cm}|p{4.8cm}|}
% % \hline
% % \textbf{ID} & \textbf{Title} & \textbf{Year} & \textbf{Summary} & \textbf{Training} & \textbf{Tasks} & \textbf{Model} & \textbf{Performance} \\
% % \hline
% % 1 & Traffic Anomaly Detection Model Using K-Means and Active Learning & 2022 & K-means + active learning for boosting model accuracy & Unsupervised & Novelty Detection & KMeans + Active Learning & ACC: 90.20\%, F1: 94.90\%, FPR: 0.01\% \\
% % \hline
% % 2 & Autoencoder-Based Unsupervised Intrusion Detection Using MSCNN-LSTM-AE & 2022 & Deep ensemble for detecting novel attacks & Unsupervised & Novelty Detection & MSCNN-LSTM-AE + Isolation Forest & AUC: 0.96–0.97 \\
% % \hline
% % 3 & Robust Unsupervised Network Intrusion Detection & 2023 & Self-supervised context reconstruction & Unsupervised & Novelty Detection & RUIDS & ACC: 92–99\% \\
% % \hline
% % 4 & RegalAnd: Lightweight NIDS for IoT Gateways & 2022 & DNN-based lightweight IDS & Supervised & Multiclass Classification & DNN & ACC: 99.93\% \\
% % \hline
% % 5 & DL Technique Using RNN & 2022 & RNN-based IDS & Supervised & Multiclass Classification & RNN/LSTM/GRU & ACC: 94–85\% \\
% % \hline
% % 6 & DCNN+BLSTM Hybrid IDS & 2022 & Hybrid CNN + BLSTM & Supervised & Multiclass Classification & CNN + BLSTM + DNN & ACC: 99\% \\
% % \hline
% % 7 & Deep Residual CNN & 2021 & Residual CNN + improved Gazelle optimizer & Supervised & Multiclass Classification & DRCNN & ACC: 91.56\%, F1: 97.14\% \\
% % \hline
% % 8 & Fed-ANIDS & 2023 & Federated + autoencoder & Semi-supervised & Novelty Detection & Autoencoder + Fed. Learning & ACC: 95.50\% \\
% % \hline
% % 9 & Unsupervised AD for IoT: Landscape & 2021 & Survey + 17 models benchmark & Semi-supervised & Novelty Detection & 17 Models & - \\
% % \hline
% % 10 & Adaptive DL IDS to Zero-Day Attacks & 2023 & Adaptable DL NIDS & Supervised & Multiclass + Novelty & SOCNN + LoF + iNNE & UDR: 97.09\% \\
% % \hline
% % 11 & Robust Detection of Unknown DoS Attacks & 2023 & Hybrid learning-based detection & Semi-supervised & Novelty Detection & - & - \\
% % \hline
% % 12 & \textbf{Ours} & N/A & OC-NFST for multi-class novelty detection & Semi-/Supervised & Multi-task & OC-NFST & AUCROC: 97.58\%, F1: 86.92\% \\
% % \hline
% % \end{tabular}%
% % }
% % \end{table*}

% \end{document}


































































































































































































































































































































































































































































































































































































































































































































































